\documentclass[11pt]{amsart}
\usepackage[T1]{fontenc}
\usepackage[utf8]{inputenc}
\usepackage{lmodern}
\usepackage{amsmath,amssymb,amsthm}
\usepackage[margin=1in]{geometry}
\usepackage{booktabs}
\usepackage{array,longtable,pdflscape}
\usepackage[expansion=false]{microtype}
\usepackage{enumitem}
\usepackage[hidelinks]{hyperref}
\newtheorem{theorem}{Theorem}
\numberwithin{theorem}{section}
\newtheorem{proposition}[theorem]{Proposition}
\newtheorem{lemma}[theorem]{Lemma}
\newtheorem{corollary}[theorem]{Corollary}
\newtheorem{claim}{Claim}[theorem]
\newtheorem{lemmaextra}{Lemma}[section]
\renewcommand{\thelemmaextra}{\thesection.\Alph{lemmaextra}}
\theoremstyle{remark}
\newtheorem{remark}[theorem]{Remark}
\newtheorem*{remark*}{Remark}
\numberwithin{equation}{section}
\providecommand{\tightlist}{\setlength{\itemsep}{0pt}\setlength{\parskip}{0pt}}
\title{The Hodge Conjecture for Odd-Degree Fermat Fourfolds}
\author{Hongmin Li}
\address{Graduate School of Frontier Sciences, The University of Tokyo}
\email{lihongmin@edu.k.u-tokyo.ac.jp}
\keywords{Hodge conjecture, Fermat fourfolds, algebraic cycles, character sums, conductor descent}
\subjclass[2020]{14C25, 14C30, 14J35}
\date{September 7, 2026}
\hypersetup{pdftitle={The Hodge Conjecture for Odd-Degree Fermat Fourfolds},pdfauthor={Hongmin Li}}

\begin{document}

\begin{abstract}

Let
\[
X_m^4=\{x_0^m+\cdots+x_5^m=0\}\subset \mathbf P^5
\]
be the Fermat fourfold of odd degree \(m\). We prove that the rational cycle
class map
\[
CH^2(X_m^4)_{\mathbf Q}\longrightarrow
H^4(X_m^4,\mathbf Q)\cap H^{2,2}(X_m^4)
\]
is surjective for every odd \(m\ge 3\). The proof is uniform in the conductor
and reduces every odd conductor to a fixed finite list of endpoint families,
which are then realized by explicit cycles. We encode primitive rational Hodge
structures by unit orbits of length-six Fermat characters, resolve their exact
orders by divisor-lattice Fourier projectors, and prove a descent-or-packet
classification. Every terminal packet is realized by an algebraic
correspondence built from Fermat-surface divisors, Aoki standard cycles, Ran
joins, inverse-pair contractions, and finite power maps. Same-degree surface
completion resolves active high-three and high-five orders without imposing
divisibility on their complements. A sharp global \(\ell^1\) ledger then
controls ambient one-small completions and the four-face system at the
mixed base.

\end{abstract}
\maketitle
\section{Introduction}\label{sec:1}

The Hodge conjecture for Fermat varieties admits an unusually explicit
arithmetic formulation. Primitive cohomology decomposes into one-dimensional
eigenspaces for the diagonal automorphism group \(G_m\). A character line
\(V_m(A)\) has Hodge type \((2,2)\) precisely when \(A\) has age three;
the full rational block \(V_m[A]_{\mathbf Q}\) has this type precisely when
every unit translate of \(A\) has age three. What is difficult is not
detecting Hodge classes but
realizing every rational Galois block by algebraic cycles. Decomposable
characters are controlled by lower-dimensional Fermat varieties, while
standard characters are represented by Aoki's cycles. Composite conductors
also contain nonstandard characters for which neither construction applies
directly.

Let \(e_{[A]}\in\mathbf Q[G_m]\) be the rational projector onto
\(V_m[A]_{\mathbf Q}\). If a rational-coefficient algebraic cycle \(Z\)
has nonzero projection onto \(V_m(A)\), then
\(w=e_{[A]}\operatorname{cl}(Z)\) is a nonzero rational vector in
\(V_m[A]_{\mathbf Q}\). Since this block is a multiplicity-one irreducible
rational representation, \(\mathbf Q[G_m]w=V_m[A]_{\mathbf Q}\).
The projector and deck translates are algebraic correspondences, so they
give algebraic generators of the whole block, as explained in Section 3.

Our result is the following.

\begin{theorem}\label{theorem:1.1}

For every odd integer \(m\ge3\), the cycle class map
\[
\operatorname{cl}\colon CH^2(X_m^4)_{\mathbf Q}\longrightarrow
H^4(X_m^4,\mathbf Q)\cap H^{2,2}(X_m^4)
\]
is surjective.

Consequently every odd-degree Fermat fourfold satisfies the rational Hodge
conjecture in all codimensions.
\end{theorem}

The last assertion follows because codimensions zero and four are immediate,
codimension one is the Lefschetz \((1,1)\) theorem, and codimension three follows
from hard Lefschetz applied to divisor classes. Thus codimension two is the only
substantive case.

\subsection{Relation to earlier work}\label{sec:1.1}

Shioda \cite{shioda1979}, Shioda--Katsura \cite[Theorem 1.7, pp.~101--102]{shiodakatsura1979}, Ran
\cite[pp.~138--140]{ran1980}, and Aoki \cite[Theorem 2-1, p.~388]{aoki1987} developed the character
calculus and the geometric correspondences used to construct large classes of
cycles on Fermat hypersurfaces. Aoki \cite[Corollary 2-3, p.~388]{aoki1987} proves the
Hodge conjecture for Fermat varieties in every dimension whenever
\(m=p^e\) is a prime power, with no restriction on the exponent.
Aoki's later work \cite[Theorem 0.1, p.~177]{aoki2000} treats abelian varieties
isogenous to factors of powers of the Fermat-curve Jacobian \(J_m\).
The Jacobian-power conclusion is recorded in
\cite[Theorem 1.2, p.~102]{aoki2002} for prime powers, twice prime powers, and
\[
m=2^a3^b5^c7^d,\qquad c=0\ \text{or}\ d=0,
\]
where the exponents are nonnegative. These statements concern
Fermat-type abelian varieties and powers of \(J_m\). The relation to
hypersurfaces is explicit in \cite[Corollary 3.2, p.~184]{aoki2000}: algebraicity
for the Fermat hypersurfaces in all dimensions is transported to the
Fermat-type abelian varieties. The proof of \cite[Theorem 0.1, p.~185]{aoki2000}
first establishes that hypersurface input for the stated degree families.
In particular, the odd families \(3^a5^b\) and \(3^a7^b\) are already
unbounded precedents.

More recently, da Silva
\cite[Proposition 3.1, p.~4]{dasilva2021}, in the published version, proved the Hodge
conjecture for Fermat fourfolds of degree \(m\le100\) with \((m,6)=1\).
It explicitly states that condition \(P_{33}\) fails
\cite[Proposition 3.4, p.~5]{dasilva2021}. Jumagulov \cite[Theorem 1.1, p.~2]{jumagulov2026} proved the result
for every odd \(3\le m\le199\), explicitly using an exhaustive machine census
\cite[Abstract, p.~1]{jumagulov2026}. Its Theorems 1.2--1.4 also give the split,
two-pair standard, and level-33-to-66 closure constructions used below.
The contribution here is the exhaustiveness of these algebraic leaves
for arbitrary odd conductors, established by the short-support,
correction, and ambient-completion arguments. For fourfolds this removes
the restrictions on prime-factor patterns in the existing unbounded
families and covers every odd degree. The present bibliographic and
source checks were completed on 2026-09-07.

Schoen \cite[Theorem 3.4, pp.~145--146]{schoen1989} proves the generalized Hodge conjecture
for a selected Hodge substructure on certain generalized Prym varieties under
a quasi-simple branching hypothesis. This is an important precedent for
transferring algebraic classes from Fermat sources to cyclic-cover geometry,
but it does not prove the Hodge conjecture for the Fermat source itself. In
the non-quasi-simple case, \cite[Remark 3.5, p.~146]{schoen1989} instead records the Fermat
Hodge conjecture as a sufficient input.

Kang stated a stronger result: \cite[Corollary 3.2, p.~227]{kang2016} asserts the generalized
Hodge conjecture for Fermat varieties of dimensions three and four in arbitrary
degree. Its equation \((3.3)\) invokes the product estimate of
\cite[Proposition 2.2(iv), pp.~521--522]{kang2015}. We retain Kang's result as a
statement-level precedent, not an input. The proof here is independent of
refined motivic dimension and addresses the odd-degree ordinary rational
Hodge conjecture directly.

\emph{Note on the tensor formula.} Write \(\mathcal F^pH\) for the largest
rational sub-Hodge structure of \(H\) contained in the Hodge-filtration
space \(F^pH_{\mathbf C}\), and put \(\mathcal F_t=\mathcal F^{-t}\)
on homology, as in \cite[pp.~520--521]{kang2015}. Here \(H_j(X,\mathbf Q)\) is
Borel--Moore homology and \(W_{-j}\) is its pure weight-\(-j\) part.
The product proof \cite[p.~522]{kang2015} uses the equality
\[
\begin{gathered}
\mathcal F_m W_{-i}H_i(X_1\times X_2,\mathbf Q)
\\=\bigoplus_{j+\ell=i}\mathcal F_m
  \bigl(W_{-j}H_j(X_1,\mathbf Q)\otimes W_{-\ell}H_\ell(X_2,\mathbf Q)\bigr)\\
=\bigoplus_{j+\ell=i}\ \bigoplus_{t+s=m}
  \mathcal F_tW_{-j}H_j(X_1,\mathbf Q)\otimes
  \mathcal F_sW_{-\ell}H_\ell(X_2,\mathbf Q).
\end{gathered}
\]
The second equality fails in general. For a positive-genus curve \(C\),
put \(H=H^1(C,\mathbf Q)\). Then \(\mathcal F^1H=0\), whereas
the Künneth component of the diagonal is a nonzero rational \((1,1)\)
class in \(\mathcal F^1(H\otimes H)\). Thus the corresponding
factorwise sum \(\sum_{a+b=1}\mathcal F^aH\otimes\mathcal F^bH\)
is zero and misses that class; dually the printed equality fails in the
\(j=\ell=m=1\) homology component. This identifies a gap in the
product argument, without disproving either the product estimate itself
or Kang's final Hodge-conjecture statement.

\subsection{Structure of the proof}\label{sec:1.2}

Let \(A=(a_0,\ldots,a_5)\) be a nonzero Fermat character modulo \(m\).
Its rational primitive block is the full orbit of \(A\) under
\(U_m=(\mathbf Z/m\mathbf Z)^\times\). The Hodge condition is
\[
\sum_{i=0}^5\langle ta_i\rangle_m=3m\qquad(t\in U_m).
\tag{1.1}\label{eq:1.1}
\]
An inverse pair leaves a Hodge quadruple and is realized by a
Fermat-surface divisor, a join, and inverse-pair contraction. We may
therefore assume that \(A\) is inverse-pair-free.

The proof uses three arithmetic steps.

\begin{enumerate}[label=\arabic*.,leftmargin=*]\tightlist
\item Exact-order Fourier rows and the sharp signed ledger isolate maximal
active orders. The vector-valued terminal-kernel classification of
Section 2 and the small matrices of Section 6 retain the actual
parity and equality cases.
\item Theorem 7.1 resolves every active order divisible by nine, allowing
arbitrary five-depth. Theorem 7.3 then resolves every active order
divisible by twenty-five by an explicit correction estimate.
Corollary 7.4 either gives a same-degree algebraic certificate or
removes common content, reducing both small-prime exponents to at
most one.
\item At a maximal large part, Theorem 7.5 completes the ambient one-small
faces with energies four, five, or six. The mixed base is treated
by the four-face system of Section 8 and its zero-top branch in
Section 10. The finite endpoints and their geometry are given in
Section 9. Every omitted-factor row is used only after the ledger
is exhausted or an actual ambient completion has been obtained.
\end{enumerate}

Same-degree surface completion in Lemma 3.3 allows the complementary
surface to remain at degree \(m\). Every leaf has an actual Chow
certificate, and every change of degree is transported by the finite
power-map identity \eqref{eq:2}. Section 10 assembles the proof without further
normal-form or uniqueness assumptions.

\subsection{Machine-checked finite certificates}\label{sec:1.3}

The accompanying Lean 4 sources (Lean 4.32.0, Mathlib) concern finite
integer predicates: the cyclic additive matrices of Appendix A, the
pure one-small exception levels of Appendix B, and the explicit
level-66 identities used in Section 9.2. In particular,
\(\texttt{Mixed35Direct.essential\_energy}\) verifies the
augmented-energy bound for the enumerated essential representatives;
the pure-small certificates identify the complete canonical sets at
levels 21, 33, and 39, and the level-66 certificates check the Hodge
rows and occurrence identities of the exceptional bridge. The Lean
sources, Python comparison scripts, and axiom reports are supplied in
the \(\texttt{certificates/}\) directory. The certificate proofs use
the standard Lean axioms \(\texttt{propext}\),
\(\texttt{Classical.choice}\), and \(\texttt{Quot.sound}\), without
\(\texttt{sorry}\) or \(\texttt{native\_decide}\).

Appendix A is an independent finite certificate. The essential-cross
exclusion of Lemma 8.2 is proved analytically through row differences
and column Fourier modes, without interpreting the augmented energy
as actual mass. The retained finite statements do not formalize the
new terminal definition or the infinite dependency chain.

Theorems 1.1, 4.1, and 10.2, the arbitrary-conductor arguments of
Sections 6--10, and the Chow realization are ordinary mathematical
proofs. Neither the geometric inputs nor the passage from the unit
equations to actual cycles is supplied as a Lean theorem. The finite
certificates verify their stated finite objects; their success alone
does not imply the arbitrary-conductor Hodge theorem.

\section{Fermat characters and rational Hodge blocks}\label{sec:2}

Let
\[
G_m=(\mu_m)^6/\mu_m^{\mathrm{diag}}
\]
act diagonally on \(X_m^4\). Write a character as
\[
A=(a_0,\ldots,a_5)\in (\mathbf Z/m\mathbf Z)^6,
\qquad \sum_i a_i=0,
\]
with each \(a_i\ne0\). The standard Fermat character decomposition and age
formula \cite[p.~385]{aoki1987} give
\[
H^4(X_m^4,\mathbf C)_{\mathrm{prim}}
=\bigoplus_A V_m(A),
\tag{2.1}\label{eq:2.1}
\]
where every nonzero summand \(V_m(A)\) is one-dimensional. If
\(\langle a\rangle_m\in\{1,\ldots,m-1\}\) denotes the standard representative,
then the middle-Hodge condition needed here is
\[
V_m(A)\subset H^{2,2}
\quad\Longleftrightarrow\quad
\sum_i\langle a_i\rangle_m=3m.
\tag{2.2}\label{eq:2.2}
\]

The rational block associated with \(A\) is
\[
V_m[A]_{\mathbf Q}\otimes_{\mathbf Q}\mathbf C
=\bigoplus_{t\in U_m/\operatorname{Stab}(A)}V_m(tA).
\tag{2.3}\label{eq:2.3}
\]
It is a rational Hodge structure of type \((2,2)\) exactly when \eqref{eq:1.1} holds.

More generally, write
\[
X_m^d=\{x_0^m+\cdots+x_{d+1}^m=0\}\subset\mathbf P^{d+1}.
\]
A Hodge tuple of even length \(L\) at level \(m\) is a multiset of
nonzero residues satisfying
\[
\sum_{a\in A}\langle ta\rangle_m=Lm/2\qquad(t\in U_m).
\tag{1}\label{eq:1}
\]
Its grade is \(L/2\). We write \(\operatorname{Alg}_m(A)\) for
algebraicity of its complete rational primitive block on \(X_m^{L-2}\).
All varieties and cycles are over \(\mathbf C\); rationality refers to
cycle coefficients, not to fields of definition of individual subvarieties.

For a subgroup \(J\) of a unit group, let \(E_J\) denote averaging under
its multiplicative action. For \(p^e\parallel q\), put
\[
J_{p,e}=\begin{cases}
U_p,&e=1,\\
\ker(U_{p^e}\to U_{p^{e-1}}),&e>1,
\end{cases}
\qquad
P_q=\prod_{p^e\parallel q}(I-E_{J_{p,e}}).
\]
We set \(P_1=I\). A terminal fibre is a coset of one of the groups
\(J_{p,e}\), with every other primary coordinate fixed. The characters
surviving \(P_q\) are precisely those of conductor \(q\).

We use \(\alpha_3\) for the nonprincipal character of \(U_3\),
normalized by \(\alpha_3(1)=1\) and \(\alpha_3(2)=-1\).
In a small-coordinate Fourier mode, \(\alpha_5\) denotes a
nonprincipal character of \(U_5\); the principal small character
is written as \(1\).

\begin{lemma}[centered tensor spark]\label{lemma:2.1}

For finite sets \(X_i\) of sizes \(n_i\ge2\), put
\(c_i(x)=\delta_x-n_i^{-1}\mathbf1\). The columns
\(\bigotimes_i c_i(x_i)\) admit no nonzero relation on fewer than
\(\min_i n_i\) distinct sites. The statement also holds for vector
coefficients.

\end{lemma}
\begin{proof}

The only relation among all centered columns of one factor is their sum.
In a proposed relation on fewer than \(\min_i n_i\) sites, fewer than
\(n_1\) first labels occur, so their columns are independent. Grouping by
that label and inducting on the remaining factors proves the assertion.
Apply linear functionals to treat vector coefficients.
\end{proof}

\begin{lemma}[twelve-site classification for large factors]\label{lemma:2.2}

Let all primes dividing \(r>1\) be at least seven and let \(V\) be a
characteristic-zero vector space. If \(0\ne h:U_r\to V\), \(P_rh=0\), and
\(|\operatorname{supp}h|\le12\), then \(h\) is constant and nonzero on one
complete terminal fibre, or on each of two parallel complete \(U_7\)-fibres.
The second possibility has twelve sites and requires \(7^1\parallel r\).

\end{lemma}
\begin{proof}

Work first on one product of terminal cosets, sorting its sizes
\(n_1\le n_2\le\cdots\). With at least two factors, \(n_1\ge6\) and
\(n_2\ge10\). Put \(P'=\prod_{i>1}(I-E_i)\). The kernel equation says that
\(P'h_a\) is independent of the first coordinate \(a\). If one slice is
zero, all slices are in \(\ker P'\), and each nonzero slice has at least
\(n_2\) sites. There is at most one; induction gives a complete fibre.
If every slice is nonzero and a singleton slice exists, the difference
between it and any other slice has at most \(14-n_1\le8<n_2\) sites.
Lemma 2.1 makes every slice identical to that singleton. If no singleton
exists, \(2n_1\le12\) forces \(n_1=6\) and two sites in every slice.
Slice differences have at most four sites and vanish, giving two parallel
six-fibres. One factor alone has the constant kernel. Finally different
terminal-coset blocks are preserved by \(P_r\); two nonzero blocks under
the bound must each be a six-fibre. The only terminal group of size six
is \(U_7\).
\end{proof}

\begin{corollary}[odd large-factor kernel]\label{corollary:2.3}

A nonzero integral odd \(h\in\ker P_r\) of norm at most twelve is
\[
h=\epsilon(\mathbf1_{U_7\times\{z\}}-
\mathbf1_{U_7\times\{-z\}}),
\qquad r=7R,\quad R>1,\quad (R,7)=1,\quad\epsilon=\pm1.
\]
In particular its norm is twelve. A nonzero kernel with support at most
six is a single \(U_7\)-fibre.

\end{corollary}
\begin{proof}

A fibre is equal to its negative only when it has no other fixed primary
coordinate and its varying subgroup contains \(-1\). This is precisely
the full group \(U_p\) at a prime level; a constant there is even.
Otherwise oddness requires the distinct negative fibre. Apply Lemma 2.2,
then integrality and the norm bound.
\end{proof}

\begin{lemma}[a five-factor under a six-site bound]\label{lemma:2.4}

On a product with one factor of size five and all other terminal sizes at
least six, a nonzero centered-product kernel function on at most six sites
is constant on a complete five-fibre or on a complete six-fibre. The same
holds for vector coefficients.

\end{lemma}
\begin{proof}

Group by the five-coordinate. If a slice is zero, every other slice is a
residual kernel relation of support at least six, so there is exactly one;
its shortest relation is a full six-fibre by the proof of Lemma 2.2.
If all five slices are nonzero, one is a singleton and every slice
difference has at most three sites. All slices equal that singleton by
residual spark. The one-factor case is immediate.
\end{proof}

We will repeatedly use the following explicit detector. For \(Q=p^e\),
\(p\ge11\), primitive characters of prescribed even parity with exponents
\(2,4\), or prescribed odd parity with exponents \(1,3\), take distinct
values at seven. Equality would imply \(7^2=1\pmod Q\), impossible because
\(p\nmid48\). Both choices are primitive, including at higher depth.
On a product of such factors, fix all but one character and use these two
choices to avoid one specified product value at seven.

\begin{lemma}[short odd total variation]\label{lemma:2.5}

Let a permutation \(T\) commute with a fixed-point-free involution \(N\),
and suppose every \(T\)-orbit has length at least three. Let \(f\) be
integral and satisfy \(f(Nx)=-f(x)\), and set
\[
k=\|f\|_1,\qquad V(f)=\sum_x|f(Tx)-f(x)|.
\]
If \(k=2\), then \(V(f)=4\). If \(k=4\), then \(V(f)\ge4\), and equality
means that \(f\) consists of two consecutive unit values and their negatives,
on one negative-stable cycle or a pair of exchanged cycles. If \(k=6\) and
\(V(f)=0\), the support is a pair of exchanged three-cycles with opposite
constant unit values.

\end{lemma}
\begin{proof}

On any cycle the variation is at least \(2(\max f-\min f)\). On a
nonnegative cycle with zeros and maximum \(a>0\) it is at least \(2a\).
Variations add over cycles. A negative-stable cycle has even length and
\(N\) acts by its half-period; a nonzero odd function on it has both signs
and variation at least four. Exchanged cycles have equal norms and variations.

For \(k=2\), \(f\) is a single unit antipodal delta. Its two sites and their
translates are distinct, since \(Tx=\pm x\) would give an orbit of length
at most two. Hence the variation is four. For \(k=4\), twice a single
antipodal delta has variation eight. Otherwise the positive and negative
parts each have two unit entries. An exchanged pair has norm two on each
cycle; because its length is at least three, variation two is possible
only when the two entries have one sign and form a single consecutive
plateau. Any second pair contributes at least four more variation. On a
negative-stable cycle, variation four permits just one positive and one
negative plateau: an additional positive or negative run would add two
transitions. The two plateaux are interchanged by \(N\), so their two
sites are consecutive. This proves the equality classification. Finally,
zero variation means constancy on every cycle. Oddness forces zero on
negative-stable cycles. An exchanged pair of nonzero cycles costs
\(2\ell|a|\) with \(\ell\ge3\), so total norm six forces \(\ell=3\) and
\(|a|=1\).
\end{proof}

The characterwise form of Theorem 1.1 is the following.

\begin{proposition}\label{proposition:2.6}

If \(A\) is a nonzero length-six character modulo an odd integer \(m\) and
\[
\sum_i\langle ta_i\rangle_m=3m
\qquad(t\in U_m),
\]
then the complete rational block \(V_m[A]_{\mathbf Q}\) is generated by
codimension-two algebraic cycles on \(X_m^4\).

\end{proposition}

\section{The typed Chow calculus}\label{sec:3}

We use character operations only together with their geometric source, target,
codimension, and field of definition. The classical inputs are the
Lefschetz \((1,1)\) theorem, Ran's algebraic join and its primitive
injectivity, Aoki's join and inverse-pair contraction, and Aoki's
standard-five cycles. The rational block is always the complete unit orbit.

\begin{lemma}[join, contraction, and finite descent]\label{lemma:3.1}

Let \(T_i\) be middle Fermat characters of length \(2d_i\) and grade \(d_i\).

\begin{enumerate}[label=\arabic*.,leftmargin=*]\tightlist
\item If the rational blocks of \(T_1\) and \(T_2\) are algebraic, then the block
of their ordered juxtaposition \(T_1*T_2\) is algebraic on
\(X_m^{2(d_1+d_2)-2}\). The induced correspondence adds one Tate factor.
\item If \(\delta=(u,-u)\) is a nonzero inverse pair and the block of
\(T*\delta\) is algebraic, then the block of \(T\) is algebraic. Repeated
inverse pairs may be contracted successively.
\item If \(n\mid m\), \(c=m/n\), and \(T=cT_0\) has length \(2d\), then
algebraicity of the block of \(T_0\) on \(X_n^{2d-2}\) is equivalent
to algebraicity of the block of \(T\) on \(X_m^{2d-2}\).
\end{enumerate}

\end{lemma}
\begin{proof}

Ran's joining-line incidence correspondence gives the first implication:
\cite[Proposition 4.6 and Corollary 4.7, p.~140]{ran1980} give injectivity on primitive
homology and identify its action on character lines with ordered
juxtaposition. The second implication follows directly from Aoki
\cite[Theorem 1-4(ii), p.~388]{aoki1987}, applied with the nonnegative even parameter
\(s=0\) and the nonzero inverse-pair character \(\delta\). Thus the
length-two contraction is included in the stated classical theorem.
Ran's proof also starts with two zero-dimensional Fermat hypersurfaces.

For the third assertion use
\[
\pi_c\colon X_m^{2d-2}\longrightarrow X_n^{2d-2},
\qquad [z_i]\longmapsto[z_i^c].
\]
Then
\[
\pi_c^*V_n(T_0)=V_m(cT_0),
\qquad \pi_{c,*}\pi_c^*=c^{2d-1}\operatorname{id}.
\]
The map has degree \(c^{2d-1}\). Its quotient group acts trivially on
the block of \(cT_0\), so \(\pi_c^*\pi_{c,*}\) is also that nonzero
scalar on this block. Hence
\[
\operatorname{Alg}_m(cT_0)\quad\Longleftrightarrow\quad
\operatorname{Alg}_n(T_0).
\tag{2}\label{eq:2}
\]
Surjective unit reduction makes both maps compatible with the complete
rational orbit.
\end{proof}

The rational group-algebra projector onto a unit orbit is an algebraic
rational linear combination of deck correspondences. A nonzero rational
vector in this multiplicity-one rational irreducible representation
generates it under the deck group. Thus a cycle with nonzero projection
on one complex character line supplies the entire rational block after
rational projection and deck translates. No averaging over arithmetic
Galois conjugates is required to assert nonvanishing.

\begin{lemma}[standard packets]\label{lemma:3.2}

Let \(m=5d\) be odd with \(d>1\), put \(p=5\), and let \(S_p(a)\) be the length-six Fermat character
\[
S_p(a)=(a,a+d,\ldots,a+(p-1)d,-pa),
\qquad a\not\equiv0\pmod d.
\]
The rational block of \(S_p(a)\) is algebraic. This includes nonprimitive
parameters: with \(g=\gcd(a,d)\), the tuple is the common-content
inflation by \(g\) of the primitive-parameter standard tuple at degree
\(m/g=5(d/g)\).

\end{lemma}
\begin{proof}

This is Aoki's standard cycle \cite[Theorem 2-1, p.~388]{aoki1987}. The standard
element range is \(d/\gcd(a,d)>2\) \cite[p.~387]{aoki1987}, and the construction
assumes \(\gcd(a,d)=1\). For a nonprimitive parameter with
\(a\not\equiv0\pmod d\), divide \(g=\gcd(a,d)\). The reduced
\(d/g\) is odd and greater than one, hence at least three. Apply Aoki's
theorem there and return by \eqref{eq:2}.
\end{proof}

\begin{lemma}[two same-degree completion rules]\label{lemma:3.3}

Let \(A\) be a Hodge sextuple, and write
\[
L_p(c)=\{c+jm/p:0\le j<p\}\pmod m.
\]

\begin{enumerate}[label=\arabic*.,leftmargin=*]\tightlist
\item If \(3\mid m\), \(c\not\equiv0\pmod{m/3}\), and
\(A=L_3(c)\uplus R\), then \(\operatorname{Alg}_m(A)\) holds.
\item If \(5\mid m\), \(c\not\equiv0\pmod{m/5}\), \(d\in L_5(c)\), and
\(A=(L_5(c)\setminus\{d\})\uplus\{y,z\}\), then
\(\operatorname{Alg}_m(A)\) holds. No divisibility of \(y,z\) is assumed.
\end{enumerate}

The second rule gives a uniform fibre criterion for the two-pair standard
closure of Jumagulov \cite[Theorem 1.3, p.~4]{jumagulov2026}: the two pairs and the
standard sextuple plus surface quadruple are exactly those in \eqref{eq:5}.
The criterion places the complementary surface at the ambient degree.
The new exhaustiveness arguments below determine when these established
geometric closures must occur at an arbitrary odd conductor.

\end{lemma}
\begin{proof}

The elementary fibre identity is
\[
\sum_{x\in L_p(c)}\langle tx\rangle_m
=\frac{p-1}{2}m+\langle ptc\rangle_m.
\tag{3}\label{eq:3}
\]
It follows by reducing to the \(p\) equally spaced lifts of \(tc\) modulo
\(m/p\). For the first assertion put
\[
S=L_3(c)\uplus\{-3c\},\qquad B=R\uplus\{3c\}.
\]
Both are nonzero Hodge quadruples by \eqref{eq:3} and \eqref{eq:1}. The literal identity
\[
A\uplus(3c,-3c)=S\uplus B
\tag{4}\label{eq:4}
\]
proves the claim by joining their surface divisor classes and contracting
the displayed pair. The join first reaches codimension three on \(X_m^6\);
the contraction returns to codimension two on \(X_m^4\). For the second
assertion set
\[
S=L_5(c)\uplus\{-5c\},\qquad B=(5c,y,z,-d).
\]
Here \(S\) is standard-five, while \eqref{eq:1} and \eqref{eq:3} give grade two for \(B\).
All entries are nonzero, and
\[
A\uplus(d,-d)\uplus(5c,-5c)=S\uplus B.
\tag{5}\label{eq:5}
\]
Join the standard cycle and the divisor to codimension four on \(X_m^8\),
then contract twice to codimension two on \(X_m^4\). If \(c\) is not
primitive, divide \(\gcd(c,m/5)\) before applying Aoki and use \eqref{eq:2} to
return. The hypothesis \(c\not\equiv0\pmod{m/5}\) makes the reduced
standard block nonzero.
\end{proof}

\begin{remark}[same-degree terminal certificates]\label{remark:3.4}

The surface \(B\) in either rule is at the same level \(m\). The terminal
condition in Section 4 therefore allows these actual surface/standard
join-and-contraction certificates at degree \(m\), without requiring the
surface factor to be inflated from \(m/3\) or \(m/5\). No divisibility
condition on the complementary entries is part of either completion rule.
\end{remark}

\section{The conductor-uniform classification}\label{sec:4}

Call a Hodge sextuple \emph{decomposable} if it contains an inverse pair.
A nondecomposable sextuple is called \emph{terminal} when it has one of the
following actual Chow certificates: joins of same-degree Fermat-surface
divisor blocks followed by inverse-pair contractions; an Aoki
standard-five block, optionally joined with a same-degree surface divisor
and contracted; a split block obtained from a divisor on a product of
Fermat curves with complementary character types; or the exceptional
level-33 construction of Section 9.2. The definition includes unit and
permutation translates and finite power-map transports of these
certificates. In particular, both completion rules of Lemma 3.3 are
terminal certificates at degree \(m\); no prior inflation of their
surface factor from \(m/3\) or \(m/5\) is required.

\begin{theorem}[descent or packet]\label{theorem:4.1}

Let \(m\ge3\) be odd and let \(A\) be a nonzero length-six character satisfying
the unit equations \eqref{eq:1.1}. Then one of the following holds.

\begin{enumerate}[label=\arabic*.,leftmargin=*]\tightlist
\item \(A\) is decomposable;
\item \(A\) is terminal;
\item there exist a proper divisor \(n\mid m\), \(c=m/n>1\), and a Hodge sextuple
\(A_0\) modulo \(n\) such that \(A=cA_0\).
\end{enumerate}

Every terminal case in the second alternative has an actual Chow realization. In particular, the
alternatives are stable under passage to the complete rational unit orbit.
\end{theorem}

The proof of Theorem 4.1 is given in Section 10 after the high-depth,
ambient one-small, and mixed-base arguments. Proposition 5.2 records
the well-founded degree changes. Assuming the classification for the
moment, the main theorem follows.

\begin{proof}[Proof of Theorem 1.1 from Theorem 4.1]

Apply Theorem 4.1 to a rational Hodge block. A decomposable character is
algebraic by Lemma 3.1 and the Lefschetz \((1,1)\) theorem on the Fermat-surface
factor. A terminal character is algebraic by Lemmas 3.1--3.3 and the explicit
packet constructions. In the remaining case, the finite power-map identity \eqref{eq:2} reduces the claim to
a proper divisor of the conductor. Strong induction on \(m\) proves
Proposition 2.6 and hence the surjectivity of the primitive cycle class map.
The nonprimitive part of \(H^4\) is generated by the square of the hyperplane
class.
\end{proof}

\section{Exhaustive branch structure and termination}\label{sec:5}

The conductor bookkeeping separates the arithmetic arguments while
retaining the full ambient Hodge sextuple. Maximality isolates rows
whose conductor contains the selected large part; actual completion
or a sharp ledger equality controls the other supports.

\begin{proposition}[branch coverage]\label{proposition:5.1}

Every odd conductor \(m\ge3\) belongs to exactly one of the following
regimes.

\begin{enumerate}[label=\arabic*.,leftmargin=*]\tightlist
\item If \(15\nmid m\), either all prime divisors are at least seven,
or a unique \(s\in\{3,5\}\) divides \(m\). In the latter case
\[
m=s^eM,\qquad(M,s)=1,\qquad p\ge7\text{ for every }p\mid M.
\tag{5.1}\label{eq:5.1}
\]
The possibilities \(M=1,e=1\), \(M=1,e\ge2\),
\(M>1,e=1\), and \(M>1,e\ge2\) are disjoint and exhaustive.
\item If \(15\mid m\), then uniquely
\[
m=3^a5^bM,\qquad a,b\ge1,\qquad(M,15)=1.
\tag{5.2}\label{eq:5.2}
\]
The cases \((a,b)=(1,1)\), \(a\ge2,b=1\),
\(a=1,b\ge2\), and \(a,b\ge2\) are disjoint and exhaustive.
\end{enumerate}

\end{proposition}
\begin{proof}

This is unique factorization, with three and five distinguished.
Theorem 7.1 applies to active three-depth at least two in any of these
regimes, and Theorem 7.3 subsequently handles active five-depth at
least two. The partition introduces no compatibility assumption
between two primary walls.
\end{proof}

\begin{proposition}[well-founded descent]\label{proposition:5.2}

Every nonterminal change of degree used in Theorem 4.1 removes a common
content \(c>1\) and therefore strictly lowers the positive conductor.
In particular, for \(m=3^a5^bM\), Corollary 7.4 reduces to
\[
3^{\min(a,1)}5^{\min(b,1)}M
\tag{5.3}\label{eq:5.3}
\]
unless a same-degree leaf has already been obtained. At a reduced
level, if no nontrivial large part is active, all entries descend to
degree three, five, or fifteen. Any such passage from a larger degree
is again strict. Consequently no reduction cycle exists.

\end{proposition}
\begin{proof}

If \(A=cA_0\) and \(c>1\), then \(m/c<m\). Surjective unit reduction
preserves the Hodge equations, and \eqref{eq:2} transports algebraicity in both
directions. The positive integers are well ordered.
\end{proof}

\begin{proposition}[maximal rows and incomparable supports]\label{proposition:5.3}

At a level with three- and five-adic exponents at most one, choose a
divisibility-maximal active large part \(r\). A Fourier row with large
conductor \(r\) receives contributions only from large parts divisible
by \(r\). Maximality removes proper multiples, but does not itself
remove incomparable supports.

The local classification gives a contradiction, an actual terminal
certificate, or a grade-two core whose two ambient complements form
an inverse pair. In the one-small sector, energy four gives that
core, energy five gives the forced level-33 completion, and energy
six exhausts the sextuple. Thus globalization follows from actual
completion or ledger exhaustion.

\end{proposition}
\begin{proof}

Equation \eqref{eq:7} proves the restriction on contributing orders. Theorem
7.5 proves the ambient one-small completions. Lemmas 8.1--8.4 resolve
the nonzero mixed face, and Section 10 treats the two-small and empty
sectors. Each use of an omitted-factor row there follows an explicit
twelve-unit equality; the lower-energy cases instead complete the
actual Hodge tuple before excluding external supports.
\end{proof}

\section{Exact-order Fourier rows and small matrices}\label{sec:6}

An inverse pair in a Hodge sextuple leaves a Hodge quadruple, so that case
is algebraic by Section 3. We henceforth assume inverse-pair-freeness.
For an exact order \(q\mid m\), define
\[
\mu_q(x)=\#\{a\in A:a=(m/q)x\},\qquad
F_q(x)=\mu_q(x)-\mu_q(-x),\qquad x\in U_q.
\]
At odd \(m\) the sharp ledger is
\[
F_q(-x)=-F_q(x),\qquad
\sum_{q\mid m,\ q>1}\|F_q\|_1=12.
\tag{6}\label{eq:6}
\]
Signed vanishing implies actual emptiness, because an inverse pair is
forbidden. Use the nonconjugate Fourier convention
\[
\widehat f(\chi)=\sum_x f(x)\chi(x),\qquad
B_q(x)=\langle x\rangle_q/q-\tfrac12,\qquad
C_q(\chi)=\sum_{x\in U_q}B_q(x)\chi(x).
\]
Then \(T_sf(y)=f(sy)\) has multiplier \(\overline{\chi(s)}\); in
particular its support is the inverse image under multiplication by \(s\).
If \(\chi\) is a primitive odd character of conductor \(f\), the
standard Dirichlet identities give
\[
C_f(\chi)=-L(0,\chi)\ne0.
\]
Indeed \(L(0,\chi)=-i\tau(\chi)L(1,\bar\chi)/\pi\), with
\(L(1,\bar\chi)\ne0\) and \(|\tau(\chi)|^2=f\); see
\cite[equations 25.15.5, 25.15.9, 25.15.10, and 27.10.11]{dlmf}. The exact-order
equation is
\[
\sum_{f\mid q\mid m}\frac{C_q(\bar\chi)}{\varphi(q)}
\widehat F_q(\chi)=0.
\tag{7}\label{eq:7}
\]
To derive it, rewrite \eqref{eq:1} as \(\sum_{q,x}F_q(x)B_q(tx)=0\), and take
its normalized Fourier coefficient against \(\bar\chi(t)\). An
order-\(q\) term vanishes unless \(f\mid q\); otherwise the substitution
\(u=tx\) gives \(\chi(x)C_q(\bar\chi)/\varphi(q)\). Unit reduction to
any divisor is surjective. For \(p\nmid f\), induction of a nonprincipal
character satisfies
\[
C_{pf}(\chi)=(1-\chi(p))C_f(\chi).
\tag{8}\label{eq:8}
\]
This follows by subtracting the multiples of \(p\) from
\(\sum_{a=1}^{pf}a\chi(a)\) and dividing by \(pf\); the sum of
\(\chi\) over a period is zero. Thus at a maximal active exact order,
\eqref{eq:7} and oddness give \(P_qF_q=0\). A row only sees orders divisible by
its conductor. Incomparable supports are not discarded until a ledger
equality proves their emptiness or an actual ambient completion resolves
the sextuple.

\begin{lemma}[small matrix classification]\label{lemma:6.1}

An integral additive \(3\times4\) matrix of norm at most six consists of
complete three-columns, a complete four-row, an essential mass-five matrix,
or an essential mass-six matrix. For \(3\times5\), the alternatives are
complete three-columns, a complete five-row, or an essential mass-six
matrix. In the latter cases the signed coefficient sums are respectively
\(\pm1,\pm2\) for \(3\times4\), and \(\pm2\) for \(3\times5\).
An integral additive \(2\times5\) matrix of norm at most six is
row-invariant of even mass, or has mass five with exactly one unit entry
per column.

\end{lemma}
\begin{proof}

Write entries as \(u_i+v_j\), choosing integer potentials. In the
three-row case the norm of each column is at least \(\max u_i-\min u_i\),
so with four or five columns that integral range is at most one. Range
zero gives complete three-columns. Otherwise, up to a gauge, row
permutation, and overall sign, use \(u=(1,0,0)\). A column is
\((t+1,t,t)\). The norm bound permits only \(t=0,-1\), since any other
column has cost at least four, making the total at least seven with four
columns. With \(n\) columns and \(b\) columns of the second kind, the
mass is \(n+b\). Thus \(b\le2\) for \(n=4\) and \(b\le1\) for
\(n=5\). The coefficient sums follow at once: \(3-2=1\),
\(2-4=-2\), and \(4-2=2\). A row-invariant matrix of mass six may
contain a three-column with multiplicity two, which is not omitted.
For two rows, the row difference \(d\) has \(5|d|\le6\). If \(d=0\),
the mass is even. If \(d=\pm1\), each column has odd norm, so the total
is odd and at least five; under the bound it equals five, with one unit
entry in every column.
\end{proof}

\begin{lemma}[residual boundary with row-lift freedom]\label{lemma:6.2}

Let a signed multiplicity restricted to an oriented terminal block have
norm at most six, and suppose its small centered projector followed by the
large residual projector vanishes. Either every residual slice satisfies
the small kernel equation, or all six sites form a residual \(U_7\)-fibre
and the projected small coefficient vector is constant on that fibre.
If all small factor sizes are at least three, all small labels and signed
coefficients are constant. For small sizes \(2,5\), the five-label is
fixed and the quantity \(\sigma\alpha_3(a)\) is constant; the two-row
label \(a\) itself need not be fixed.

\end{lemma}
\begin{proof}

If there is no residual factor, its projector is the identity and the
small equation holds immediately. Otherwise apply the residual projector
to the vector-valued small projection. Fewer than six residual sites give
zero by Lemma 2.1; exactly six give a constant-vector complete
\(U_7\)-fibre. Equality in the original integral norm leaves one unit
coefficient at each site. Centered delta columns for a set of size at
least three are pairwise nonproportional. Equality of nonzero simple
tensors consequently fixes all labels and coefficients in that case.
On a two-set the two centered columns are negatives; multiplying by its
nontrivial character records precisely that sign freedom. This gives the
stated \(2,5\) alternative rather than an unjustified fixed-row assertion.
\end{proof}

\section{High primary depth and ambient one-small completion}\label{sec:7}

\subsection{Direct high-three closure}\label{sec:7.1}

\begin{theorem}[any active three-depth at least two is algebraic]\label{theorem:7.1}

Let \(m\) be odd and \(A\) an inverse-pair-free Hodge sextuple. If some
entry has exact additive order divisible by nine, then \(A\) contains a
complete three-fibre, a standard-five sextuple, or four members of a
nonzero five-fibre. In particular \(\operatorname{Alg}_m(A)\) follows
from Lemma 3.3 and Aoki's standard cycle.

\end{theorem}
\begin{proof}

Choose a divisibility-maximal active order \(q\) divisible by nine.
Every active proper multiple would again be divisible by nine, so its
primitive row is isolated. Write \(q=3^i5^jr\), where \(i\ge2\) and
\((r,15)=1\); \(j\) may be zero. The terminal three-kernel has size
three and odd order, so a terminal block and its negative are disjoint.
Each oriented block has integral norm at most six.

Apply Lemma 6.2. A residual boundary exhausts all six physical entries,
fixes their three-primary label up to a common sign, and forces all
entries to have exact order \(q\). Divide the common content \(m/q\).
Modulo \(3^i\) their sum is \(\pm6z\), with \(z\) a unit. This is
nonzero because \(i\ge2\). Thus no such boundary occurs.

All residual slices therefore satisfy the small equation. If \(j=0\),
a nonzero slice is constant on a three-fibre. If \(j=1\), use the
\(3\times4\) case of Lemma 6.1; if \(j\ge2\), use \(3\times5\).
A complete signed three-column gives an actual complete three-fibre
(choose its positive or negative orientation), and is algebraic by
Lemma 3.3. A five-row at depth \(j\ge2\) gives five actual members of
\(L_5(c)\); the sixth is \(-5c\) by the zero-sum congruence. A four-row
when \(j=1\) gives four actual members of \(L_5(c)\) and is covered by
the same-degree completion rule, regardless of its two complements.

An essential mass-six slice exhausts the ledger. All its three-primary
labels are congruent to one unit modulo three, but its signed coefficient
sum is \(\pm2\). The six actual entries, after dividing common content,
therefore have nonzero sum modulo three, a contradiction.

Only the essential mass-five \(3\times4\) slice remains. Its five
effective labels have sum \(S\) that is a unit modulo every prime dividing
\(q\). At three the coefficient sum is \(\pm1\). At five, if \(b_0\)
is the exceptional column, the weighted sum is \(\pm3b_0\). At each
residual prime the common label is multiplied by \(\pm1\). Hence the
sixth actual entry is \(-(m/q)S\) and also has order \(q\). Reapply the
just-established six-mass classification to this full stratum. If the
sixth entry lay in a new independent small slice, that slice would contain
only one unit entry, contrary to the minimum-support bound for the
centered-product kernel in Lemma 2.1. A
three-fibre or incomplete five-fibre is an algebraic leaf; the essential
six-mass or residual boundary is impossible. Thus every alternative gives
an asserted leaf or contradiction.
\end{proof}

\begin{corollary}\label{corollary:7.2}

Unless \(A\) is already algebraic by one of these explicit certificates,
every active order has three-adic valuation at most one. Hence, if
\(3^a\parallel m\), all coordinates are divisible by
\(3^{\max(a-1,0)}\).
\end{corollary}

This conclusion allows arbitrarily large ambient five-depth. It eliminates
every active high-three order before the high-five argument is applied.

\subsection{High-five closure and the correction estimate}\label{sec:7.2}

\begin{theorem}[high-five closure after high-three removal]\label{theorem:7.3}

Let \(m\) be odd and let \(A\) be an inverse-pair-free Hodge sextuple
with no active order divisible by nine. If some active order is divisible
by twenty-five, then \(\operatorname{Alg}_m(A)\) has a standard-five or
surface-join-contraction certificate.

\end{theorem}
\begin{proof}

Choose a maximal active base \(h=q/3^{v_3(q)}\) among orders divisible
by twenty-five. Write \(h=5^jr\), \(j\ge2\), \((r,15)=1\). Only the
orders \(h\) and \(3h\) can contribute to a conductor-\(h\) row. Set
\(D=F_h\) and \(F=F_{3h}\), defining the latter to be zero if
\(3h\nmid m\). If \(F\ne0\), its primitive row gives the
\(U_3\times J_5\) small equation with the residual projector.

First eliminate the residual boundary in Lemma 6.2. It uses all six
entries. Let \(\eta_3\) be the odd character of \(U_3\), choose a
primitive even character \(\theta\) at five-depth \(j\), and primitive
even characters on every other remaining primary factor. Twisting
\(\theta\) by the quadratic character modulo five preserves both
evenness and full depth, and changes its value at seven by \(-1\).
Thus one of the two resulting odd characters of conductor \(3h/7\) has
a nonzero omitted-seven Euler factor. Its coefficient of \(F\) is
nonzero: on the positive fibre the product \(\sigma\eta_3(a)\) is fixed
by Lemma 6.2, and the negative fibre gives the same contribution. The
isolated row \eqref{eq:7} contradicts this.

Now Lemma 6.1 applies on every small slice. A complete five-row is a
standard-five leaf. An essential mass-five slice has five actual effective
labels whose sum is a unit at every primary factor: its three-coordinate
weighted sum is \(\pm5u\), and its signed coefficient sum on the five
and residual factors is \(\pm1\) or \(\pm3\). The sixth entry therefore
has the same exact order \(3h\). A six-mass additive \(2\times5\) slice
is row-invariant; another noninvariant slice would already cost five, and
a separate nonzero slice costs at least two. Thus, after the standard
leaves and contradictions, the entire top face is
\[
F(a,y)=v(y),\qquad v(-y)=-v(y),\qquad k=\|v\|_1\in\{2,4,6\}.
\]
Its full signed norm is \(2k\). The edge row gives
\[
PC=0,\qquad C=D-(T_3v-v),\qquad
P=(I-E_{J_5})\prod_{p^e\parallel r}(I-E_{J_{p,e}}).
\tag{9}\label{eq:9}
\]
We now prove \(C=0\) by a norm estimate.

Put \(g=T_3v-v\). The ledger gives \(\|D\|_1\le12-2k\) and
\(\|g\|_1\le2k\), so \(\|C\|_1\le12\). All terminal blocks are
exchanged in pairs under negation because \(J_5\) has odd order. On an
oriented half, \(C\) has norm at most six. Lemma 2.4, integrality, and
the budget show that a nonzero correction must have the exact form
\[
C=\epsilon(\mathbf1_L-\mathbf1_{-L}),\qquad
\epsilon=\pm1,\qquad |L|=\rho\in\{5,6\}.
\tag{10}\label{eq:10}
\]
There is just one such antipodal block pair. A six-fibre is in an
exponent-one \(U_7\)-direction. In either case \(L\) has a fixed nonzero
five-primary residue \(z\) modulo \(5^{j-1}\).

Let \(\Lambda=L\cup(-L)\). A source pair \(\{x,-x\}\) of \(v\)
contributes to \(g\) at \(x,-x,x/3,-x/3\). At most one of its two
antipodal pairs can lie in \(\Lambda\), since otherwise
\(3z=\pm z\pmod{5^{j-1}}\). Consequently
\[
\|g|_\Lambda\|_1\le k,\qquad
\|D\|_1=\|C+g\|_1\ge2\rho-k.
\tag{11}\label{eq:11}
\]
For \(\rho=6\), this contradicts \(\|D\|_1\le12-2k\) for every
\(k>0\). For \(\rho=5\) it contradicts that bound when \(k\ge4\).
When \(\rho=5,k=2\), \(v\) is a single unit antipodal delta. If
\(g\) meets \(\Lambda\), its other antipodal pair lies outside
\(\Lambda\) and contributes two unavoidable units. Hence
\(\|D\|_1\ge10-2+2=10>8\); if there is no meeting the bound is even
larger. This excludes \eqref{eq:10} and proves \(C=0\).

Thus
\[
D=T_3v-v,\qquad k+\tfrac12\|T_3v-v\|_1\le6.
\tag{12}\label{eq:12}
\]
Every multiplication-by-three orbit in \(U_h\) has length at least
twenty, since \(5^j\mid h\): the order at five is four, and
\(v_5(3^4-1)=1\) gives \(\operatorname{ord}_{5^j}(3)=4\cdot5^{j-1}\)
by binomial lifting. These are point orbits in \(U_h\). For \(k=2\),
\(v=\pm(\delta_c-\delta_{-c})\) and the variation is four. The two top
atoms and two forced lower atoms form a Hodge quadruple; the other two
entries have grade one and are opposite. For \(k=4\), Lemma 2.5 gives
variation at least four; equality is forced. On exchanged orbits, norm
two and variation two on one orbit give two adjacent unit values; on a
negative-stable orbit the same follows from its positive and negative
plateaux. Therefore, up to translation and sign,
\[
v=\delta_c+\delta_{3c}-\delta_{-c}-\delta_{-3c}.
\]
For precision about the actual labels, put \(d=m/(3h)\) and let \(P_c\)
be the two unit lifts of \(c\in U_h\) to level \(3h\), multiplied by
\(d\). Then
\[
S_c=P_c\uplus\{3d(c/3),-3dc\}=L_3(dc)\uplus\{-3dc\}
\]
is a Hodge quadruple. The two \(S\)-blocks satisfy
\[
S_c\uplus S_{3c}=A\uplus(3dc,-3dc).
\]
This realizes the six-entry equality. For \(k=6\) the variation is zero;
an odd invariant function vanishes on negative-stable orbits, and a
nonzero pair of exchanged orbits would have norm at least forty. Hence
this case is impossible.

It remains to justify the case \(F=0\). Then \(PD=0\). An oriented
nonzero block has at most six units, so Lemma 2.4 gives a complete
five-fibre (standard-five) or six \(U_7\)-points with fixed five-primary
unit label. The latter exhausts all entries and has sum
\(\pm6z\ne0\pmod{5^j}\) after removal of common content. Thus it is
impossible, completing the proof.
\end{proof}

\begin{corollary}[simultaneous primary-depth reduction]\label{corollary:7.4}

For odd \(m=3^a5^bM\), \((M,15)=1\), a Hodge sextuple is algebraic by
an explicit same-level leaf certificate, or all its entries are divisible by
\[
c=3^{\max(a-1,0)}5^{\max(b-1,0)}.
\]
In the latter case \(A=cA_0\) is a Hodge sextuple at level
\[
m/c=3^{\min(a,1)}5^{\min(b,1)}M.
\]

\end{corollary}
\begin{proof}

Remove inverse pairs first. Theorem 7.1 eliminates any active
nine-divisible order unless a leaf is obtained. Theorem 7.3 then does the
same for twenty-five. The remaining valuation bounds are exactly the
common-content assertion. Dividing preserves all unit equations by
surjective unit reduction; use \eqref{eq:2} for algebraicity.
\end{proof}

\subsection{Ambient one-small completion}\label{sec:7.3}

All primes of \(r>1\) are at least seven and \(s\in\{3,5\}\).
Suppose at a maximal active large part only the exact-order faces
\(sr\) and \(r\) are present. This applies both to the pure one-small
setting and to the mixed base with its other two faces zero. Write
\[
F(a,y)=F_{sr}(a,y),\qquad D(y)=F_r(y),\qquad
e=\tfrac12(\|F\|_1+\|D\|_1),\qquad F\ne0.
\]
The top and edge rows are isolated by maximality; rows omitting a large
factor are not yet isolated. Formula \eqref{eq:7} proves their independence of
the ambient modulus. In unnormalized form the common rescaling from
level \(sr\) to \(m\) is \([m/(sr)]\varphi(m)/\varphi(sr)\).

\begin{theorem}[ambient completion]\label{theorem:7.5}

Under these hypotheses and the six-entry Hodge ledger, one obtains an
actual energy-four Hodge quadruple whose complements are inverse, the
energy-five prefix of a level-33 exceptional packet whose last entry is
forced, or an energy-six algebraic packet exhausting all entries.

\end{theorem}
\begin{proof}

First suppose \(r\) is not a prime. Apply \(\pi_s=I-E_{U_s}\) to the
small-coordinate vector and set \(H(y)=\pi_sF(\cdot,y)\). Then
\(P_rH=0\) and \(H(-y)=-N_sH(y)\), where \(N_s\) negates the small
coordinate. A nonzero \(H\) under twelve units, by Lemma 2.2, must use
two opposite \(U_7\)-fibres and exhaust \(\|F\|_1=12\). A
self-opposite single fibre would require \(r\) prime. For \(s=3\) the
projected coefficient \(\sigma\alpha_3(a)\) is constant along a fibre;
for \(s=5\) the map \((\sigma,a)\mapsto\pi_5(\sigma\delta_a)\) is
injective, so its small label and sign are constant. Choose an odd
nontrivial \(\alpha_s\) detecting this coefficient and a primitive even
residual character on \(r/7\) avoiding the one omitted-seven value. The
detector following Lemma 2.4 provides it. After the ledger exhaustion
the conductor-\(sr/7\) row has one nonzero term and nonzero Euler factor.
Contradiction. Thus
\[
F(a,y)=b(y),\qquad b\text{ integral and odd}.
\]
The edge row gives \(P_rC=0\), where \(C=D-(T_s-1)b\). If \(s=5\),
\(k=\|b\|_1\ge2\) and
\[
\|C\|_1\le12-(s-3)k\le8,
\]
so Corollary 2.3 forces \(C=0\). If \(s=3\), the only possible
nonzero correction has norm twelve and is a pair of opposite
\(U_7\)-fibres. Equality in
\[
\|C\|_1\le\|D\|_1+\|T_3b-b\|_1\le12-2k+2k=12
\]
forces \(D,T_3b-b\) to have disjoint support inside the unit-coefficient
support of \(C\), and forces \(b,T_3b\) to have disjoint support there
too. Both supports would have residual labels \(z,-z\) on \(r/7\),
requiring \(3z=\pm z\). All primes of \(r/7\) are at least eleven, so
this is impossible. Hence also \(C=0\).

The energy is now
\[
e=\frac{s-1}{2}k+\frac12\|T_sb-b\|_1\le6.
\tag{13}\label{eq:13}
\]
For \(s=5\), necessarily \(k=2\), the variation is four, and the
complete packet is standard-five. For \(s=3\), \(k=2\) gives a
four-entry three-standard surface block. For \(k=4\), Lemma 2.5 makes
variation four equivalent to two consecutive unit values and their
negatives, so \(e=6\) and the two-surface identity from the high-five
proof applies. For \(k=6\), invariance would require
\(\operatorname{ord}_r(3)\le3\) on a nonzero odd orbit. Orders one or
two are impossible; order three forces \(r\mid26\), hence \(r=13\),
contrary to the current nonprime case. This completes the nonprime
argument.

Now let \(r=p\) be prime. The exact top kernel gives, in the zero-mean
gauge,
\[
F(a,y)=u(y)+b(a),\qquad u(-y)=-u(y),\qquad b(-a)=-b(a).
\tag{14}\label{eq:14}
\]
Integrality makes the potentials either both integral or both strict
half-integral: all differences within each potential are integral and
oddness forces twice the common fractional part to be integral. The edge
equation gives
\[
D=T_su-u
\]
exactly, since the only \(P_p\)-kernel is constant and this correction
is odd.

For \(s=5\), write \(b=(c_1,c_2,-c_1,-c_2)\). Pairing opposite small
coordinates gives
\[
\sum_a|u(y)+b(a)|=2\max(|u(y)|,|c_1|)+2\max(|u(y)|,|c_2|).
\]
If the potentials are integral, nonzero noninvariant \(u\) has norm at
least two and variation at least four, giving \(e\ge6\). Nonzero
invariant odd \(u\) costs more: \(\operatorname{ord}_p(5)\ge3\).
When \(u=0\), nonzero \(b\) already gives \(e\ge p-1\ge6\). In
the half-integral case the same lower bound \(e\ge p-1\) holds. Thus
\(e=6\) and every other stratum is empty. The remaining small-only row
consequently says \((T_p-1)b=0\). An integral nonzero \(b\) would
force \(p=7\), equality, \(u=0\), and support on a single antipodal
small pair; multiplication by \(7\equiv2\pmod5\) does not preserve it.
In the half-integral case equality is possible only at \(p=7\), but
would require an odd nowhere-zero \(u\) invariant under five, although
\(5^3=-1\pmod7\). Hence \(b=0\). Then \(u\) is an integral single
antipodal delta and the packet is standard-five.

For \(s=3\) put \(b=(c,-c)\) and \(V(u)=\|T_3u-u\|_1\). The
exact energy is
\[
e=\sum_{y\in U_p}\max(|u(y)|,|c|)+\tfrac12V(u)\le6.
\tag{15}\label{eq:15}
\]
In the integral case, \(c\ne0\) forces \(p=7\), \(|c|=1\), and
equality. Then \(V(u)=0\) and \(3^3=-1\pmod7\) give \(u=0\).
If \(c=0\), the even norm \(k=\|u\|_1\) is \(2,4\), or \(6\).
Norm two gives an actual four-entry surface block; norm four has variation
at least four, and equality gives the quasi packet. Norm six requires
three-invariance, hence \(p=13\) as above, with energy six.

In the strict half-integral case \(p\le13\), so
\(p\in\{7,11,13\}\). Also \(|c|=1/2\). At \(p=11,13\), a value
of \(|u|\ge3/2\) already increases the first term by at least two and
exceeds the budget. At \(p=7\), such a value would give first term at
least five and variation at least six, also impossible. Thus every
\(|u(y)|=1/2\). At \(p=13\) the first term is already six. At
\(p=11\) it is five, and a nonconstant sign pattern on either five-cycle
has total variation at least four on the negative pair of cycles; thus
\(u\) is three-invariant and \(e=5\). At \(p=7\), \(U_7\) is a
three-generated cycle with negation its half-period. Its three
half-period signs have one or three sign changes on the path ending at
the negative of the first; the full variations are two or six and the
energies are four or six. This proves that only \(e=4,5,6\) occur.

We finally justify actual ambient completions. At energy four with
\(c=0\) the actual core is the explicit three-standard quartet. At
energy four in the half-integral \(p=7\) case, the core has zero full
21- and 7-conductor rows by \eqref{eq:14} and \(D=T_3u-u\). Its remaining odd
row has conductor three and vanishes because \(1-\alpha_3(7)=0\).
Fourier inversion therefore proves that these four actual entries form
a Hodge quadruple. Inflating to the ambient level preserves this
assertion. Subtracting its Hodge equations from the sextuple leaves a
grade-one pair, necessarily opposite.

At energy five, \(p=11\). The five positive entries form one small row
times a coset of \(\langle3\rangle\subset U_{11}\). Their sum is
\(11\) or \(22\) modulo \(33\). Thus the ambient zero-sum equation
forces the sixth entry to have order three, and the completed tuple is a
unit/permutation translate of
\[
(m/33)(1,4,16,22,25,31).
\]
At energy six, the two faces exhaust all six entries. All actual entries
are multiples of \(m/(3p)\), so the full effective Hodge equations hold
at \(3p\). For \(p>13\) the preceding integral classification already
gave the quasi family. The remaining exact levels \(21,33,39\) are the
finite lists in Section 9, and their blocks are algebraic by the
constructions there. No omitted-large-factor row has been isolated before
one of these actual completions.
\end{proof}

\section{The exact four-face system at level \(15M\)}\label{sec:8}

Assume \(m=15M\), \((M,15)=1\). At a maximal active large part
\(r\mid M\) write its four faces as
\[
F_{35}(a,b,y),\qquad F_3(a,y),\qquad F_5(b,y),\qquad D(y),
\qquad a\in U_3,\ b\in U_5,\ y\in U_r.
\]
Here the subscripts abbreviate the exact orders \(15r,3r,5r,r\).
We write \(E_s=E_{U_s}\), and \(N_s\) for negation of the
\(U_s\)-coordinate. For a primitive large character \(\beta\), the top
equation is
\[
P_r(I-E_3)(I-E_5)F_{35}=0.
\tag{16}\label{eq:16}
\]
The edge equations, obtained from \eqref{eq:7}, are
\[
4\widehat F_3(\alpha_3,\beta)
+(1+\overline{\beta(5)})\widehat F_{35}(\alpha_3,1,\beta)=0,
\tag{17}\label{eq:17}
\]
\[
2\widehat F_5(\alpha_5,\beta)
+(1-\overline{\alpha_5(3)\beta(3)})
\widehat F_{35}(1,\alpha_5,\beta)=0.
\tag{18}\label{eq:18}
\]
The trivial-small equation is
\[
\begin{aligned}
0={}&8\widehat D(\beta)
+4(1-\overline{\beta(3)})\widehat F_3(1,\beta)
+2(1-\overline{\beta(5)})\widehat F_5(1,\beta)\\
&+(1-\overline{\beta(3)})(1-\overline{\beta(5)})
\widehat F_{35}(1,1,\beta).
\end{aligned}
\tag{19}\label{eq:19}
\]
Parity ensures the even modes cause no additional conditions.

\begin{lemma}[Pointwise additive reduction]\label{lemma:8.1}

\(F_{35}(a,b,y)=u(a,y)+v(b,y)\) with integral potentials.

\end{lemma}
\begin{proof}

Set \(D'_b(y)=F_{35}(1,b,y)-F_{35}(2,b,y)\) and
\(H(y)=\pi_5D'(y)\), where \(\pi_5=I-E_5\). Then
\(P_rH=0\), \(D'(-y)=N_5D'(y)\), and \(H(-y)=N_5H(y)\).
Thus \(\operatorname{supp}H\) is negative-stable, but no parity is imposed
on an individual rectangular difference. If \(r=1\), \(H=0\) immediately.

Suppose \(H\ne0\). Since \(|\operatorname{supp}H|\le\|F_{35}\|_1\le12\),
Lemma 2.2 applies. A single negative-stable fibre requires \(r=p\) and
\(H\) constant. Then \(N_5H=H\), so every integral representative
\(D'(y)\) is constant on antipodal small pairs and has norm at least two.
Consequently \(\|F_{35}\|_1\ge2(p-1)\); equality forces \(p=7\), norm
twelve, and
\[
D'(y)\in\{\epsilon\mathbf1_{\{1,4\}},-\epsilon\mathbf1_{\{2,3\}}\},
\qquad D'(-y)=D'(y).
\]
If two fibres occur, they are opposite \(U_7\)-fibres in \(r=7R\),
\(R>1\), and already require all twelve units of \(F_{35}\). There is
one unit entry over each residual site. Injectivity of
\((\epsilon,b)\mapsto\pi_5(\epsilon\delta_b)\) then gives
\[
D'(u,z)=\epsilon\delta_b,\qquad
D'(u,-z)=\epsilon\delta_{-b}\quad(u\in U_7).
\]
These descriptions retain the possible row choice of each original
\(F_{35}\) atom; the row difference, not a fictitious fixed row, is used
below.

Either alternative exhausts \eqref{eq:6}, so all other actual strata vanish. The
omitted-seven row is now isolated. In the lone \(r=7\) case take
\(\chi=\alpha_3\lambda_5\), where \(\lambda_5\) is quadratic. The row
difference gives \(\widehat F_{35}(\chi)=12\epsilon\ne0\), while
\(\chi(7)=-1\). In the paired case choose a primitive odd \(\beta\)
on \(R\) and the two odd quartic characters \(\gamma_\pm\) of
\(U_5\). Then \(\chi_\pm=\alpha_3\gamma_\pm\beta\) are odd, both
detect the full multiplicity with coefficient
\(12\epsilon\gamma_\pm(b)\beta(z)\), and have values
\(\pm i\beta(7)\) at seven. At least one value is not one. In either
case \eqref{eq:8} and primitive odd nonvanishing contradict \eqref{eq:7}. Thus \(H=0\).
Fixing one row and one column supplies integral additive potentials directly.
\end{proof}

\begin{lemma}[essential-cross exclusion]\label{lemma:8.2}

For \(r>1\), no additive slice of \(F_{35}\) genuinely varies in both
the row and column directions.

\end{lemma}
\begin{proof}

At an oriented residual pair \(\{y,-y\}\) the small matrix has physical
mass at most six. Its integral common row difference \(d\) satisfies
\(\|F_y\|_1\ge4|d|\) and the norm is even. Thus an essential slice has
\(d=\pm1\) and mass four or six. A mass-two slice is row-invariant.

Define \(G_3(y)=F_3(1,y)-F_3(2,y)\). Equation \eqref{eq:17} gives
\[
P_r(G_3-h_3)=0,\qquad h_3=-(I+T_5)d.
\]
The essential slice and its negative give the same unit value of \(d\)
at \(y,-y\), so \(h_3\) is a nonzero constant \(c\) on the four
distinct sites \(\{\pm y,\pm y/5\}\). If the total top mass is six,
\(F_3=0\); the only possible additional mass-two slice does not change
\(d\). A four-site kernel relation contradicts Lemma 2.1.

Otherwise the top mass is four and at most two physical atoms remain.
The correction \(K=G_3-h_3\) is even and has support at most eight.
Lemma 2.2 and parity give \(K=0\) or \(r=7\) with \(K\) constant.
If \(K=0\), the four-site \(G_3\) consumes both remaining atoms. Choose
a nonprincipal column character \(\alpha_5\) detecting the nonconstant
column potential \(w\) at \(y\). Its forced edge mode is
\[
h_5=\overline{\alpha_5(3)}T_3w-w,
\]
nonzero on four distinct residual sites. With no \(F_5\) atom remaining,
this contradicts the scalar spark.

If \(r=7\) and \(K=\lambda\mathbf1\ne0\), support at most four for
\(G_3\) forces \(\lambda=-c\). Then \(G_3\) is supported on the
complementary antipodal pair and consumes at least one physical atom,
leaving at most one for \(F_5\). A single-atom \(F_5\) mode has at
most two sites. It cannot differ from \(h_5\) by a constant on \(U_7\):
a nonzero constant would require the four existing values of \(h_5\)
all to be equal, since its two absent sites would already consume the two
allowed sites. For the quadratic character \(h_5\) is odd and not
constant; for a quartic character its values at \(y\) and \(y/3\) have
ratio \(-\overline{\alpha_5(3)}=\pm i\). A zero constant leaves four
sites. Contradiction.
\end{proof}

Appendix A records the independent finite classification of these
additive matrices and its augmented-energy inequality. Lemma 8.2 uses
the row-difference and column Fourier modes above; it does not interpret
augmented energy as actual physical mass.

If any row-only slice occurs, it has exactly four atoms. An additional
mass-two column-only slice would make the top mass six and contradict
the same \(G_3\) argument. Hence the four atoms form a complete
\(U_5\)-row, that is, four members of \(L_5(c)\). Lemma 3.3 resolves
it directly. The case \(r=1\) descends to the finite level fifteen in
Section 9. We are left with the global column-only sector
\[
F_{35}(a,b,y)=v(b,y),\qquad k=\|v\|_1>0,\qquad\|F_{35}\|_1=2k.
\tag{20}\label{eq:20}
\]

\subsection{Correction lifting and small-constant closure}\label{sec:8.1}

Put
\[
(S_3v)(b,y)=v(3b,3y),\qquad g=S_3v-v,\qquad C=F_5-g.
\]
Equation \eqref{eq:18} gives \(P_r(I-E_5)C=0\), and
\[
\|C\|_1\le\|F_5\|_1+2k\le12-\|F_3\|_1-\|D\|_1.
\tag{21}\label{eq:21}
\]

\begin{lemma}[correction is small-coordinate constant]\label{lemma:8.3}

In \eqref{eq:20}, under the full Hodge system \eqref{eq:16}--\eqref{eq:19} and ledger \eqref{eq:6},
\(C(b,y)=c(y)\), with integral odd \(c\).

\end{lemma}
\begin{proof}

For \(r=1\) this is immediate. Let \(H(y)=\pi_5C(\cdot,y)\); now its
actual parity is \(H(-y)=-N_5H(y)\). If \(H\ne0\), Lemma 2.2 applies
using \(\|C\|_1\le12\). In a paired-fibre case, equality leaves one
unit coefficient of \(C\) per residual site. The injective projected
delta map fixes one small column on the first \(U_7\)-fibre and its
negative column on the opposite fibre. Thus \(C\) uses just one
antipodal small pair. A single self-opposite fibre requires \(r=p\).
Then \(C(b,y)=a(b)+f(y)\) in the odd zero-mean gauge. The potentials
are both integral or both half-integral, and
\[
\|C\|_1\ge(p-1)\|a\|_1\ge2(p-1)
\]
when \(a\ne0\). Hence \(p=7\) and every bound is an equality.
In the integral case \(a\) is a unit antipodal delta and \(f=0\),
again using only two small columns. In the half-integral case every
potential value has absolute value \(1/2\), and
\[
C=\epsilon(\mathbf1_{R\times Z}-\mathbf1_{-R\times(-Z)}),
\]
where \(R\subset U_5\) and \(Z\subset U_7\) are transversals of the
antipodal pairs, of sizes two and three.

Every case has \(\|C\|_1=12\). Equality in \eqref{eq:21} makes \(F_3=D=0\),
\(\|g\|_1=2k\), and every nonzero coefficient of \(C\) a unit.
Equality in the triangle inequality then forces \(F_5\) and \(g\) to
be disjoint, and \(v,S_3v\) to be disjoint with both supports inside
\(\operatorname{supp}C\). A two-column support cannot contain both,
since multiplication by three interchanges the two antipodal column
pairs. In the half-integral rectangle,
\(\operatorname{Agg}_5C=2\epsilon(\mathbf1_Z-\mathbf1_{-Z})\) is
nonzero and odd, where \(\operatorname{Agg}_5\) sums the four columns.
Choose an odd character \(\beta\) on \(U_7\) detecting it. Since five
generates \(U_7\), \(\beta(5)\ne1\). The equality already exhausted
all other strata. Substituting \(F_5=S_3v-v+C\) and \(F_{35}=v\) in
\eqref{eq:19} cancels the \(v\)-terms and leaves
\[
2(1-\overline{\beta(5)})\widehat{\operatorname{Agg}_5C}(\beta)=0,
\]
a contradiction. Thus \(H=0\) and the assertion follows.
\end{proof}

The top plus five-face mass satisfies \(2k+\|F_5\|_1\ge8\).
For \(k\ge4\) this is immediate. For \(k=2\), \(v\) is a unit
antipodal delta, and \(g\) has one unit in each of four distinct
residual slices. Adding a column-constant integer cannot reduce any such
slice below norm one. At most four signed units remain for \(F_3,D\);
its row difference is a \(P_r\)-kernel function with support at most
four. Hence
\[
F_3(a,y)=b_3(y)
\]
by Lemma 2.1. The trivial-small row first gives only
\[
P_rK=0,\qquad K=D-(T_3-1)b_3-(T_5-1)c.
\]
Writing \(J=2k+\|F_5\|_1\), we have
\[
4\|c\|_1\le J,\qquad J+2\|b_3\|_1+\|D\|_1\le12,
\]
so \(\|K\|_1\le12-J+2\|c\|_1\le8\). For \(c\ne0\) use
integral oddness \(\|c\|_1\ge2\) and \(J\ge4\|c\|_1\); for
\(c=0\) use \(J\ge4\). Corollary 2.3 therefore gives \(K=0\).

\begin{lemma}[small-constant closure]\label{lemma:8.4}

\(c=b_3=D=0\) in the column-only sector.

\end{lemma}
\begin{proof}

Assume \(c\ne0\). Then \(4\|c\|_1\le12\) and integral oddness give
\(\|c\|_1=2\). Since \(J\ge8\), the remaining budget allows either
\(b_3=0\), in which case \(D=(T_5-1)c\) has norm four, or
\(\|b_3\|_1=2,D=0\). In both cases \(J\le8\). Thus
\[
8=\|C\|_1\le\|F_5\|_1+\|g\|_1
\le\|F_5\|_1+2k=J\le8.
\]
As above, integrality and unit coefficients force \(v,S_3v\) to be
disjoint and supported in \(C\). Its residual support is just
\(\{z,-z\}\), so \(3z=\pm z\) would follow, impossible. Thus
\(c=0\). Now \(D=(T_3-1)b_3\). If \(b_3\ne0\), the cost
\(2\|b_3\|_1+\|(T_3-1)b_3\|_1\) is at least eight: norm two has
variation four and norm at least four suffices already. The top and
five-face cost is also at least eight, violating the ledger. Hence
\(b_3=D=0\).
\end{proof}

Finally \(k+\tfrac12\|S_3v-v\|_1\le6\). Every \(S_3\) orbit has
length divisible by four, owing to its \(U_5\) coordinate. Norm two
gives an actual surface core and an inverse complementary pair.
Lemma 2.5 turns norm four and variation four into two adjacent unit
values and their negatives, and the same two-surface contraction identity
used earlier. Norm six would require invariance; a negative-stable orbit
then contributes zero and an exchanged pair would cost at least eight.
Thus the nonzero mixed face always has an algebraic certificate.

\section{Finite endpoints and their Chow realizations}\label{sec:9}

\subsection{Complete finite orbit lists}\label{sec:9.1}

At level fifteen the inverse-pair-free unit/permutation orbits have
representatives
\[
(1,2,7,11,12,12),\qquad(1,4,7,10,10,13).
\]
The second is standard-five. With \(C(w)=L_3(w)\uplus\{-3w\}\), the
first obeys
\[
(1,2,7,11,12,12)\uplus(6,9)=C(2)\uplus C(6).
\]
For \(p=7,11,13\), put
\[
Q_p=(1,1+p,1+2p,3+p,3+2p,3p-9)\pmod{3p}.
\]
The complete inverse-pair-free orbit lists are
\[
\begin{array}{c|l}
p&\text{representatives}\\\hline
7&Q_7,\ (1,4,9,15,16,18),\ (1,4,10,13,16,19),\\
11&Q_{11},\ (1,4,16,22,25,31),\\
13&Q_{13},\ (1,7,16,22,34,37),\ (1,14,16,22,29,35).
\end{array}
\]
An exhaustive enumeration sets the sixth entry to \(3m\) minus the sum
of the first five nondecreasing positive entries below \(m\), retains
its legal nondecreasing range, excludes inverse pairs, and tests every
unit equation. Sorting every unit translate gives canonical orbits.
Comparison of the complete canonical sets gives the displayed lists.
The unquotiented counts at \(15,21,33,39\) are \(6,12,14,24\); their
orbit counts are \(2,3,2,3\). Appendix B retains the exact finite
certificate at \(21,33,39\).

\subsection{Chow realization of the finite leaves}\label{sec:9.2}

The \(Q_p\) blocks satisfy
\[
Q_p\uplus(3,-3)=C(1)\uplus C(3)
\]
and are algebraic. The split closure is the one in Jumagulov
\cite[Theorem 1.2, p.~4]{jumagulov2026}. The four split representatives admit
\[
\begin{array}{c|l}
21&(1,4,16)*(9,15,18),\quad(1,4,16)*(10,13,19),\\
39&(1,16,22)*(7,34,37),\quad(1,16,22)*(14,29,35).
\end{array}
\]
Each triple sums to zero at its level. Since every unit translate of
the juxtaposition is Hodge, the two curve characters have complementary
types; their diagonal rational orbit is a \((1,1)\)-block on the product
of curves. Lefschetz supplies a divisor there and Ran's algebraic join
\cite[Section 4, pp.~138--139]{ran1980} maps it nontrivially to the fourfold block
\cite[Proposition 4.6 and Corollary 4.7, p.~140]{ran1980}. This uses the divisor
on the product; algebraicity of either curve character alone is not
asserted. The rational projector argument in Section 3 gives the
complete orbit.

For the exceptional level-33 tuple
\(w=(1,4,16,22,25,31)\), we use the certificate of Jumagulov
\cite[Theorem 1.4, p.~5]{jumagulov2026}. The following \(w,a,Q,S,\beta,\gamma,\delta_0,\delta_1\)
are the same tuples as there, and the level-66 completion, contractions,
and descent by the squaring map are the same construction. Work at level
66 with \(a=2w\) and
\[
\begin{aligned}
a&=(2,8,32,44,50,62),& Q&=(1,25,44,62),\\
S&=(2,8,32,41,50,65),&\delta_0&=(1,65,25,41),\\
\beta&=(2,32,33,65),&\gamma&=(8,33,41,50),
\end{aligned}
\qquad\delta_1=(33,33).
\]
The three quadruples \(Q,\beta,\gamma\) are Hodge and
\[
a\uplus\delta_0=Q\uplus S,\qquad
S\uplus\delta_1=\beta\uplus\gamma.
\]
Join \(\beta,\gamma\) to codimension three on \(X_{66}^6\) and
contract \(\delta_1\) to obtain the codimension-two block of \(S\).
Then join with \(Q\) to codimension four on \(X_{66}^8\) and contract
the two pairs in \(\delta_0\), returning to codimension two on
\(X_{66}^4\). The self-pair \((33,33)\) is a legitimate nonzero
inverse pair at even level 66. This uses only that explicit even
auxiliary variety. The squaring map \(X_{66}^4\to X_{33}^4\) has
quotient group of order 32 acting trivially on \(a\); hence
\(\pi^*\pi_*=32\) on that block and its algebraic image at level 33
is nonzero. Reduction \(U_{66}\simeq U_{33}\) preserves the full orbit.
The finite certificate checks all displayed Hodge conditions and multiset
identities exactly.

\section{The zero mixed face and global assembly}\label{sec:10}

Retain the notation of Section 8 and first suppose \(r>1\) and
\(F_{35}=0\). If exactly one of \(F_3,F_5\) is nonzero, Theorem 7.5
applies. If both are zero, \(P_rD=0\); an odd nonzero large-factor
kernel would have norm twelve by Corollary 2.3. This exhausts all strata
and its omitted-seven row is detected by a primitive odd character on
\(r/7\) avoiding one value at seven. Thus it is impossible.

Suppose both small faces are nonzero. For \(r\) not prime, put
\(H_s=\pi_sF_s\). Its negative-stable support is strictly smaller than
twelve, since the other nonzero odd face already costs at least two
units. Lemma 2.2 gives \(H_s=0\): a single self-opposite fibre requires
\(r\) prime, and a paired fibre requires twelve sites. For \(r=p\)
prime, write \(F_s(a,y)=c_s(a)+b_s(y)\) with odd integral or
half-integral potentials. A nonzero \(c_5\) costs at least
\((p-1)\|c_5\|_1\ge12\) signed units, impossible with \(F_3\ne0\).
Thus \(c_5=0\) and \(F_5=b_5\) costs at least eight. A nonzero
\(c_3\) then costs at least \(p-1\ge6\) signed units, again
impossible. Consequently in every case
\[
F_3=b_3(y),\qquad F_5=b_5(y).
\]
Their signed costs are at least four and eight. Equality in the ledger
makes each \(b_s\) a single antipodal delta and empties \(D\) and every
other stratum. Only now are all residual rows isolated. Put
\[
K=(T_3-I)b_3+(T_5-I)b_5.
\]
The primitive odd rows of exact conductor \(r\), whose Bernoulli
coefficients are nonzero, give \(P_rK=0\). The function \(K\) is
integral and odd, and
\[
\|K\|_1\le2\|b_3\|_1+2\|b_5\|_1=8.
\]
By Corollary 2.3, a nonzero integral odd function in \(\ker P_r\)
has norm at least twelve. Hence \(K=0\), proving
\[
(T_3-1)b_3+(T_5-1)b_5=0.
\tag{22}\label{eq:22}
\]
On the basis \(e_y=\delta_y-\delta_{-y}\), indexed by
\(U_r/\{\pm1\}\), each summand is an edge between distinct vertices.
Equality of the two edges forces one of
\[
r\mid5-3,\qquad r\mid5+3,\qquad r\mid15-1,\qquad r\mid15+1.
\]
Every prime of \(r\) is at least seven, so \(r=7\) is the only
nontrivial possibility. The equations then normalize to
\(b_3=e_3,b_5=e_1\). The actual level-105 representative is
\[
(3,24,50,66,85,87)=(L_5(3)\setminus\{45\})\uplus\{50,85\},
\]
which is the example in Jumagulov \cite[Theorem 1.3, p.~4]{jumagulov2026} and is
algebraic by Lemma 3.3. Unit translations and inflation cover
the whole branch.

\begin{proposition}[no small prime]\label{proposition:10.1}

If every prime dividing an odd \(m\) is at least seven, every Hodge
sextuple is a union of inverse pairs.

\end{proposition}
\begin{proof}

For an inverse-pair-free input, take a maximal active exact order \(q\).
Corollary 2.3 makes the only nonzero possibility a pair of opposite
\(U_7\)-fibres at \(q=7R\), exhausting the ledger. Choose a primitive
odd character on \(R\) with value at seven different from one, using
the detector in Section 2. Its coefficient is nonzero and \eqref{eq:8} contradicts
the single surviving row. After removing inverse pairs from a general
sextuple, any nonempty remainder has even length at most four. Its odd
ledger is at most eight, immediately excluded by the same kernel
argument. Hence no remainder exists.
\end{proof}

\begin{theorem}[assembled odd-degree conclusion]\label{theorem:10.2}

The classical geometric inputs in Section 3 and the arithmetic arguments
above give
\[
\operatorname{cl}:CH^2(X_m^4)_{\mathbf Q}\twoheadrightarrow
H^4(X_m^4,\mathbf Q)\cap H^{2,2}(X_m^4)
\qquad\text{for every odd }m\ge3.
\]

\end{theorem}
\begin{proof}

By the classical character decomposition it suffices to realize each
Hodge sextuple. An inverse pair is a surface-divisor leaf. Otherwise
apply Corollary 7.4. Its explicit leaves are algebraic, or its finite
power descent gives a Hodge tuple at a level with three- and five-adic
exponents at most one. If neither small prime occurs, Proposition 10.1
applies. If exactly one \(s\) occurs, choose a maximal active large
part \(r\). For \(r>1\), a nonzero \(F_{sr}\) is covered by Theorem
7.5; if it is zero the sole \(F_r\) is excluded by the odd-kernel and
omitted-seven argument. For \(r=1\), all entries descend to the prime
level \(s\), where \(P_sF_s=0\) and oddness force \(F_s=0\),
contradicting activity. If both small primes occur, use the four-face
arguments in Sections 8 and 10. At \(r>1\) they give an algebraic leaf
or contradiction. At \(r=1\) every entry has order dividing fifteen;
remove common content and apply the finite level-fifteen classification
in Section 9. Every leaf uses actual cycles by Section 3 or Section 9;
every passage back to the original level uses \eqref{eq:2}. The nonprimitive part
is the square of the hyperplane class.
\end{proof}

\begin{proof}[Proof of Theorem 4.1]

The preceding proof gives the stated classification as well as
algebraicity. An inverse pair is the first alternative. Corollary 7.4
either supplies a same-degree terminal certificate or removes common
content; when that content exceeds one, this is the third alternative.
At the reduced level, the maximal-large-part arguments of Theorem 7.5
and Sections 8 and 10 give an actual terminal certificate or a
contradiction. If the maximal large part is one, removing common content
reduces to degree three, five, or fifteen. The first two are impossible
for an inverse-pair-free input, and degree fifteen has the two explicit
certificates of Section 9. All certificates are preserved by unit
translation, permutation, and finite power-map transport. This proves
Theorem 4.1 with the terminal definition of Section 4.
\end{proof}

\appendix
\section{The \texorpdfstring{\(2\times4\)}{2 x 4} cyclic additive certificate}\label{sec:appA}

We retain the complete finite cyclic additive certificate. Its augmented
energy is a finite matrix invariant; the analytic essential-cross
exclusion in Lemma 8.2 does not identify it with actual mass. Order the four
\(U_5\)-columns by the multiplication-by-three cycle
\[
 (1,3,4,2).
\]
Let \(M=(m_{ij})\in\mathbf Z^{2\times4}\) satisfy
\[
 m_{ij}=u_i+v_j,
 \qquad
 0<\|M\|_1\le6.
\tag{A.1}\label{eq:A.1}
\]
Equivalently, after flattening by rows, the three equations
\[
 m_{1j}-m_{10}=m_{0j}-m_{00},
 \qquad j=1,2,3,
\tag{A.2}\label{eq:A.2}
\]
hold.  The exact equivalence group is generated by the row swap, cyclic
rotation of the four columns, and multiplication of all coefficients by
\(-1\).  Reflections and arbitrary column permutations are not allowed,
because they do not preserve the multiplication-by-three dynamics.

There are two exhaustive enumerations.  The first visits the full integer
\(\ell^1\)-ball in eight coordinates:
\[
 \#\{x\in\mathbf Z^8:\|x\|_1\le6\}=40081
\]
and retains the 288 nonzero solutions of \eqref{eq:A.2}.  The second writes
\[
\begin{aligned}
(m_{00},m_{01},m_{02},m_{03})
   &=(a,a+c_1,a+c_2,a+c_3),\\
(m_{10},m_{11},m_{12},m_{13})
   &=(a+d,a+d+c_1,a+d+c_2,a+d+c_3).
\end{aligned}
\tag{A.3}\label{eq:A.3}
\]
The norm bound implies
\[
 |a|\le6,\qquad |d|,|c_j|\le12,
\]
so enumeration in these ranges is independently exhaustive.  The two
canonical class sets agree.  They contain twenty-nine classes, with norm
profile
\[
 \#\{\|M\|_1=2,4,6\}=(1,9,19).
\tag{A.4}\label{eq:A.4}
\]

A class is essential when both \(d\ne0\) and the four entries of its first
row are not all equal.  For a vector \(w\), set
\[
 q(w)=\min_{c\in\mathbf Z}\sum_i|w_i+c|.
\tag{A.5}\label{eq:A.5}
\]
Then
\[
 q_{\rm row}=|d|,
 \qquad
 q_{\rm col}=q(m_{00},m_{01},m_{02},m_{03}).
\tag{A.6}\label{eq:A.6}
\]
The eleven essential representatives and their augmented energies
\[
 E(M)=\|M\|_1+2q_{\rm row}+2q_{\rm col}
\tag{A.7}\label{eq:A.7}
\]
are as follows.  Each displayed vector is flattened by rows.

\[
\begin{array}{c|r|c|r}
M&\|M\|_1&(q_{\rm row},q_{\rm col})&E(M)\\\hline
(-1,-1,-1,0;\ 0,0,0,1) & 4 & (1,1) & 8\\
(-1,-1,0,0;\ 0,0,1,1) & 4 & (1,2) & 10\\
(-1,0,-1,0;\ 0,1,0,1) & 4 & (1,2) & 10\\
(-2,-1,-1,-1;\ -1,0,0,0) & 6 & (1,1) & 10\\
(-2,-1,-1,0;\ -1,0,0,1) & 6 & (1,2) & 12\\
(-2,-1,0,-1;\ -1,0,1,0) & 6 & (1,2) & 12\\
(-2,-1,0,0;\ -1,0,1,1) & 6 & (1,3) & 14\\
(-2,0,-1,-1;\ -1,1,0,0) & 6 & (1,2) & 12\\
(-2,0,-1,0;\ -1,1,0,1) & 6 & (1,3) & 14\\
(-2,0,0,-1;\ -1,1,1,0) & 6 & (1,3) & 14\\
(-2,0,0,0;\ -1,1,1,1) & 6 & (1,2) & 12\\
\end{array}
\]

Thus every essential class has \(E(M)\ge8>6\). This is the independent
finite inequality certified by the table. As a diagnostic, quotienting
incorrectly by the full symmetric
group on four columns leaves only sixteen classes and collapses thirteen
genuine cyclic classes.

The supplementary verifier
\( \texttt{verify\_mixed35\_c4\_additive\_classification.py} \)
uses only integer arithmetic.  It asserts equality of the two enumerations,
recomputes \eqref{eq:A.4}--\eqref{eq:A.7}, and fixes the canonical payload by SHA-256
\[
\texttt{6884de68c535501c5d08250dbb0d297fda791e05c767291166bce208600127dd}.
\]
The hash is an integrity check on the finite table, not evidence for any of
the analytic arbitrary-conductor steps in Sections 6--10.

\section{The pure one-small exception levels}\label{sec:appB}

We give the finite classification used after the analytic one-small energy
and ambient-transport arguments have reduced the pure tail to
\[
 m=3p,\qquad p\in\{7,11,13\}.
 \tag{B.1}\label{eq:B.1}
\]
The equivalence group is exactly
\[
 U_{3p}\times\mathfrak S_6.
 \tag{B.2}\label{eq:B.2}
\]
The symmetric group is removed by sorting the sextuple; multiplication by
\(-1\) is already contained in \(U_{3p}\) and is not divided out again.

Put
\[
 Q_p=(1,1+p,1+2p,3+p,3+2p,3p-9).
 \tag{B.3}\label{eq:B.3}
\]
This represents the quasi orbit. Up to \eqref{eq:B.2}, the complete lists of
inverse-pair-free grade-three Hodge sextuples are
\[
\begin{array}{c|l}
p&\text{orbit representatives}\\ \hline
7&
Q_7,\ (1,4,9,15,16,18),\ (1,4,10,13,16,19),\\
11&
Q_{11},\ (1,4,16,22,25,31),\\
13&
Q_{13},\ (1,7,16,22,34,37),
(1,14,16,22,29,35).
\end{array}
\tag{B.4}\label{eq:B.4}
\]
The second orbit at \(p=11\) is the level-33 exceptional block, whose
algebraicity certificate is due to Jumagulov \cite[Theorem 1.4, p.~5]{jumagulov2026}
(arXiv:2608.18134). The remaining
nonquasi representatives are precisely the split packets appearing in
Section 9.2.

Here is an exhaustive certificate for \eqref{eq:B.4}. For fixed \(m=3p\), sort the
six entries and enumerate the first five by combinations with repetition.
The \(t=1\) Hodge row determines the last entry uniquely:
\[
 a_6=3m-\sum_{i=1}^5a_i.
 \tag{B.5}\label{eq:B.5}
\]
Retain it only when \(a_5\le a_6<m\), discard every tuple containing an
opposite pair, and test
\[
 \sum_{i=1}^6\langle ta_i\rangle_m=3m
 \qquad(t\in U_m).
 \tag{B.6}\label{eq:B.6}
\]
Finally replace every survivor by the lexicographically least sorted unit
translate. This visits every unordered sextuple once before quotienting and
therefore proves exhaustiveness. The numbers retained before taking unit
orbits, followed by the orbit counts, are
\[
\begin{array}{c|c|c}
p&\text{retained sextuples}&\text{unit orbits}\\ \hline
7&12&3\\
11&14&2\\
13&24&3.
\end{array}
\tag{B.7}\label{eq:B.7}
\]
Direct comparison of the canonical sets gives exactly \eqref{eq:B.4}.

The supplementary verifier
\[
\begin{gathered}
\texttt{calculations/bundles/fermat-level15m-pure-small-exception-levels/}\\
\texttt{verify\_pure\_small\_exception\_levels.py}
\end{gathered}
\]
uses only integer arithmetic and asserts the canonical-set equality, not
merely the counts \eqref{eq:B.7}. On OMEN it returns
\[
\texttt{pure\_exception\_levels=PASS}.
\]

In the preceding energy reduction, energy four contains an actual
grade-two block; the two remaining entries have grade one and are forced to
be opposite, so this case returns to the decomposable branch. Energy five
occurs only at \(p=11\): the strict-half-integral prefix is unique and the
ambient sum determines the sixth entry, giving the level-33 exceptional
representative in \eqref{eq:B.4}. The energy-six survivors are the quasi and split
packets listed there. Their Chow realizations are those already used in
Section 9.2.

Appendix B certifies only the effective levels \eqref{eq:B.1}. The arbitrary-\(M\)
transport, full-depth residual-factor detectors, and four-face
globalization are the analytic arguments in the body of the paper.

\begin{landscape}

\section{Branch, budget, and geometric dependencies}\label{sec:appC}

Table \ref{tab:assembly} records the dependencies of Theorem 10.2,
separately from the finite certificate lists in Section 9 and Appendices
A--B. Write \(\delta_u=(u,-u)\), and let \(B\) denote an actual Hodge
quadruple. In the last column, a sequence such as \(2,2\to6\to4\)
records Fermat dimensions at the stated level: two surface-divisor
sources, their join, and inverse-pair contraction back to the fourfold.
The direct standard-five cycle has codimension two on \(X_m^4\);
joining it with a surface divisor and contracting twice gives
\(4,2\to8\to4\). Ran's primitive injectivity and Aoki's
Theorem 1-4(ii) give nonvanishing for these operations (Lemma 3.1),
and Aoki's Theorem 2-1 gives it for the standard cycle (Lemma 3.2).

The condition \(r=1\) concerns the maximal \emph{active} large part; it does
not assert that the ambient large part \(M\) is one. In every geometric
row, the full rational block is obtained by applying the rational deck
projector of Section 3 to a cycle with nonzero character-line projection
and taking its deck translates.

\begingroup
\footnotesize
\setlength{\tabcolsep}{3pt}
\renewcommand{\arraystretch}{1.12}
\begin{longtable}{@{}>{\raggedright\arraybackslash}p{0.12\linewidth}>{\raggedright\arraybackslash}p{0.13\linewidth}>{\raggedright\arraybackslash}p{0.18\linewidth}>{\raggedright\arraybackslash}p{0.17\linewidth}>{\raggedright\arraybackslash}p{0.19\linewidth}>{\raggedright\arraybackslash}p{\dimexpr0.21\linewidth-30pt\relax}@{}}
\caption{Branches, budgets, and geometric certificates.}\label{tab:assembly}\\
\toprule
Branch & Maximal active order & Twelve-unit budget & When omitted-factor rows are available & Actual multiset identity or outcome & Source--target; nonvanishing \\
\midrule
\endfirsthead
\multicolumn{6}{l}{\textit{Table \thetable\ continued.}}\\
\toprule
Branch & Maximal active order & Twelve-unit budget & When omitted-factor rows are available & Actual multiset identity or outcome & Source--target; nonvanishing \\
\midrule
\endhead
\midrule
\multicolumn{6}{r}{\textit{Continued on the next page.}}\\
\endfoot
\bottomrule
\endlastfoot
Inverse pair & None needed. & Resolved before the inverse-pair-free ledger \eqref{eq:6}. & Not used. & \(A=B\uplus\delta_u\). & \(0,2\to4\); Lemma 3.1.\\
\addlinespace
High three-depth & \(q=3^i5^jr\), \(i\ge2\), maximal among active orders divisible by nine. & Residual boundary or essential mass six exhausts all entries; mass five forces the sixth into order \(q\). & Not used; maximality isolates the primitive \(q\)-row. & \eqref{eq:4}, \eqref{eq:5}, or \(A=L_5(c)\uplus\{-5c\}\). & \(2,2\to6\to4\), standard cycle on \(X_m^4\), or \(4,2\to8\to4\); Lemmas 3.1--3.3, Theorem 7.1.\\
\addlinespace
High five-depth, after high three & Maximal active base \(h=5^jr\), \(j\ge2\); faces \(h,3h\). & Residual boundary exhausts twelve. After \(C=0\), \(k=4\) and variation four exhaust twelve; \(k=2\) gives a four-entry core. & Omitted seven is used only after the residual boundary exhausts the ledger. & Standard-five; \(A=B\uplus\delta_u\); or \(S_c\uplus S_{3c}\) \(=A\uplus\delta_{3dc}\). & Standard cycle, \(0,2\to4\), or \(2,2\to6\to4\); Lemmas 3.1--3.3, Theorem 7.3.\\
\addlinespace
Common content & All active orders now have three- and five-depth at most one. & No exhaustion needed; valuation bounds force common divisibility. & Not used. & \(A=cA_0\), with \(c\) from Corollary 7.4; strict descent only for \(c>1\). & \(X_m^4\to X_{m/c}^4\); scalar \(c^5\) in \eqref{eq:2}.\\
\addlinespace
No small prime; or only \(D\) at the chosen large part & Maximal active exact order \(q\), or maximal active large part \(r>1\). & The only possible nonzero odd kernel is a pair of opposite \(U_7\)-fibres, of norm twelve. & Omitted seven is isolated after equality empties every other stratum. & Inverse-pair-free case impossible; without small primes, \(A=\biguplus_{i=1}^3\delta_{u_i}\). & \(0,2\to4\) for inverse pairs; Lemma 3.1, Proposition 10.1 and Section 10.\\
\addlinespace
One small face, \(e=4\) & Maximal active large part \(r>1\); faces \(sr,r\). & Eight units form an actual Hodge quadruple; its two complements are inverse. & No omitted-large-factor row is needed for this ambient completion. & \(A=B\uplus\delta_u\). & \(0,2\to4\); Lemma 3.1, Theorem 7.5.\\
\addlinespace
One small face, \(e=5\) & \(s=3\), \(r=11\), selected as a maximal active large part. & Ten units form the prefix; the zero-sum equation forces the sixth entry, of order three. & No omitted-large-factor row is used before actual completion. & \(A=(m/33)w\), up to units and permutation; \(a\uplus\delta_0=Q\uplus S\), \(S\uplus\delta_1=\beta\uplus\gamma\). & At 66: \(2,2\to6\to4\), then \(4,2\to8\to4\); \(X_{66}^4\to X_{33}^4\) has scalar 32; Section 9.2, then \eqref{eq:2}.\\
\addlinespace
One small face, \(e=6\) & Maximal active large part \(r>1\); prime cases \(r=p\). & \(\|F_{sr}\|_1+\|F_r\|_1=12\); all other strata empty. & Effective remaining rows are now available. The earlier paired-fibre alternative uses omitted seven only after \(\|F_{sr}\|_1=12\). & Standard-five; \(S_c\uplus S_{3c}\) \(=A\uplus\delta_{3dc}\); or the \(Q_p\), split-triple, and level-33 identities of Section 9.2. & Standard cycle; \(2,2\to6\to4\); divisor on \(X_m^1\times X_m^1\to X_m^4\) via Ran \cite[Proposition 4.6, p.~140]{ran1980}; or the 66-to-33 certificate.\\
\addlinespace
Mixed \(F_{35}\ne0\) & Maximal active large part \(r>1\); faces \(15r,3r,5r,r\). & Exceptional top kernels cost twelve. Nonconstant \(C\) forces equality in \eqref{eq:21}. Surviving column norm four with variation four exhausts twelve. & Lemma 8.1: omitted seven after top exhaustion. Lemma 8.3: residual row after equality in \eqref{eq:21}. Earlier actual completions need neither. & \eqref{eq:5}; \(A=B\uplus\delta_u\); or the two-surface contraction identity of Theorem 7.3. & \(4,2\to8\to4\), \(0,2\to4\), or \(2,2\to6\to4\); Lemmas 3.1--3.3 and 8.1--8.4.\\
\addlinespace
Mixed \(F_{35}=0\), both small faces nonzero & Maximal active large part \(r>1\). & The faces cost at least four and eight; equality makes each \(b_s\) one antipodal delta and empties all other strata. & All residual rows only after equality; exact-conductor-\(r\) rows and Corollary 2.3 give \eqref{eq:22}. & \(r=7\); the level-105 representative is \((L_5(3)\setminus\{45\})\uplus\{50,85\}\), as displayed in Section 10. & \(4,2\to8\to4\) via \eqref{eq:5}; Lemmas 3.1--3.3; \cite[Theorem 1.3, p.~4]{jumagulov2026}.\\
\addlinespace
Maximal active \(r=1\) & Every active order divides \(s\), or divides 15 when both small primes occur; ambient \(M\) need not be one. & No further exhaustion; common content removes the inactive ambient large part. & Not used. & Prime levels 3 and 5 impossible; at 15, standard-five or \((1,2,7,11,12,12)\) \(\uplus(6,9)=C(2)\uplus C(6)\). & Standard cycle or \(2,2\to6\to4\) at 15; Lemmas 3.1--3.3, Section 9, and transport by \eqref{eq:2}.\\
\end{longtable}
\endgroup
\end{landscape}

\begin{thebibliography}{12}
\bibitem{aoki1987}
N. Aoki, \emph{Some new algebraic cycles on Fermat varieties}, J. Math. Soc.
   Japan \textbf{39} (1987), no. 3, 385--396.
   \href{https://doi.org/10.2969/jmsj/03930385}{doi:10.2969/jmsj/03930385}.
\bibitem{dasilva2021}
G. da Silva Jr., \emph{Notes on the Hodge conjecture for Fermat varieties},
   Experimental Results \textbf{2} (2021), e22.
   \href{https://doi.org/10.1017/exp.2021.14}{doi:10.1017/exp.2021.14}.
\bibitem{jumagulov2026}
R. Jumagulov, *The Hodge conjecture for Fermat fourfolds of odd degree at
   most 199*, arXiv:2608.18134 (2026).
\bibitem{kang2015}
S.-J. Kang, \emph{Refined motivic dimension}, Canad. Math. Bull. \textbf{58} (2015),
   519--529.
   \href{https://doi.org/10.4153/CMB-2015-018-4}{doi:10.4153/CMB-2015-018-4}.
\bibitem{kang2016}
S.-J. Kang, \emph{Refined motivic dimension of some Fermat varieties}, Bull.
   Aust. Math. Soc. \textbf{93} (2016), 223--230.
   \href{https://doi.org/10.1017/S0004972715001379}{doi:10.1017/S0004972715001379}.
\bibitem{ran1980}
Z. Ran, \emph{Cycles on Fermat hypersurfaces}, Compositio Math. \textbf{42} (1980), no. 1,
   121--142.
   \href{https://www.numdam.org/item/CM_1980__42_1_121_0/}{NUMDAM}.
\bibitem{shiodakatsura1979}
T. Shioda and T. Katsura, \emph{On Fermat varieties}, Tohoku Math. J. \eqref{eq:2}
   \textbf{31} (1979), 97--115.
   \href{https://doi.org/10.2748/TMJ/1178229881}{doi:10.2748/TMJ/1178229881}.
\bibitem{shioda1979}
T. Shioda, \emph{The Hodge conjecture for Fermat varieties}, Math. Ann. \textbf{245}
   (1979), 175--184.
   \href{https://doi.org/10.1007/BF01428804}{doi:10.1007/BF01428804}.
\bibitem{schoen1989}
C. Schoen, *Cyclic covers of \(P^v\) branched along \(v+2\) hyperplanes
   and the generalized Hodge conjecture for certain abelian varieties*, in
   \emph{Arithmetic of Complex Manifolds}, Lecture Notes in Mathematics \textbf{1399},
   Springer, 1989, 137--154.
   \href{https://doi.org/10.1007/BFb0095974}{doi:10.1007/BFb0095974}.
\bibitem{dlmf}
NIST Digital Library of Mathematical Functions, §25.15, *Dirichlet \(L\)-functions*,
    equations \href{https://dlmf.nist.gov/25.15.E5}{25.15.5},
    \href{https://dlmf.nist.gov/25.15.E9}{25.15.9}, and
    \href{https://dlmf.nist.gov/25.15.E10}{25.15.10}; see also
    \href{https://dlmf.nist.gov/27.10.E11}{27.10.11} for the primitive Gauss sum.
    \href{https://dlmf.nist.gov/25.15}{https://dlmf.nist.gov/25.15}.
    Version 1.2.7, released 15 June 2026; accessed 7 September 2026.
\bibitem{aoki2000}
N. Aoki, *Some remarks on the Hodge conjecture for abelian varieties of
    Fermat type*, Comment. Math. Univ. Sancti Pauli \textbf{49} (2000), no. 2,
    177--194.
    \href{https://doi.org/10.14992/00009819}{doi:10.14992/00009819}.
\bibitem{aoki2002}
N. Aoki, \emph{Hodge cycles on CM abelian varieties of Fermat type}, Comment.
    Math. Univ. Sancti Pauli \textbf{51} (2002), no. 1, 99--130.
    \href{https://doi.org/10.14992/00008725}{doi:10.14992/00008725}.
\end{thebibliography}
\end{document}
